% Options for packages loaded elsewhere
\PassOptionsToPackage{unicode}{hyperref}
\PassOptionsToPackage{hyphens}{url}
%
\documentclass[
  10pt,
]{article}
\usepackage{amsmath,amssymb}
\usepackage{iftex}
\ifPDFTeX
  \usepackage[T1]{fontenc}
  \usepackage[utf8]{inputenc}
  \usepackage{textcomp} % provide euro and other symbols
\else % if luatex or xetex
  \usepackage{unicode-math} % this also loads fontspec
  \defaultfontfeatures{Scale=MatchLowercase}
  \defaultfontfeatures[\rmfamily]{Ligatures=TeX,Scale=1}
\fi
\usepackage{lmodern}
\ifPDFTeX\else
  % xetex/luatex font selection
\fi
% Use upquote if available, for straight quotes in verbatim environments
\IfFileExists{upquote.sty}{\usepackage{upquote}}{}
\IfFileExists{microtype.sty}{% use microtype if available
  \usepackage[]{microtype}
  \UseMicrotypeSet[protrusion]{basicmath} % disable protrusion for tt fonts
}{}
\makeatletter
\@ifundefined{KOMAClassName}{% if non-KOMA class
  \IfFileExists{parskip.sty}{%
    \usepackage{parskip}
  }{% else
    \setlength{\parindent}{0pt}
    \setlength{\parskip}{6pt plus 2pt minus 1pt}}
}{% if KOMA class
  \KOMAoptions{parskip=half}}
\makeatother
\usepackage{xcolor}
\usepackage[margin=0.85in]{geometry}
\setlength{\emergencystretch}{3em} % prevent overfull lines
\providecommand{\tightlist}{%
  \setlength{\itemsep}{0pt}\setlength{\parskip}{0pt}}
\setcounter{secnumdepth}{-\maxdimen} % remove section numbering
\usepackage{amsmath,amssymb,amsthm}
\usepackage{longtable}
\usepackage{fvextra}
\DefineVerbatimEnvironment{verbatim}{Verbatim}{breaklines=true,fontsize=\small}
\setlength{\emergencystretch}{3em}
\ifLuaTeX
  \usepackage{selnolig}  % disable illegal ligatures
\fi
\usepackage{bookmark}
\IfFileExists{xurl.sty}{\usepackage{xurl}}{} % add URL line breaks if available
\urlstyle{same}
\hypersetup{
  pdftitle={P16.2: retained polar data on proper nonprojective varieties},
  pdfauthor={Research manuscript},
  hidelinks,
  pdfcreator={LaTeX via pandoc}}

\title{P16.2: retained polar data on proper nonprojective varieties}
\usepackage{etoolbox}
\makeatletter
\providecommand{\subtitle}[1]{% add subtitle to \maketitle
  \apptocmd{\@title}{\par {\large #1 \par}}{}{}
}
\makeatother
\subtitle{Genuine local products, intrinsic finite energy, and equations
for retained polar data}
\author{Research manuscript}
\date{5 October 2026}

\begin{document}
\maketitle

\section{Abstract}\label{abstract}

We construct a finite-energy non-Archimedean equation theory on every
smooth integral proper variety of positive dimension over an arbitrary
algebraically closed trivially valued characteristic-zero field.
Projectivity and an ample bundle on the original carrier are
unnecessary. Every specified actual big rational line bundle admits a
retained polar presentation on a smooth projective modification. Its
prescribed original-section cone consists of actual metrics on the
original carrier, with normalized compactness and genuine envelopes.
Common actual ample collars construct intrinsic local products,
independent of admissible local nets, collars and frames, including
bounded discontinuous endpoints. Relative reference floors give the
genuine nonpluripolar operator and raw mass \((A^n)\). Its intrinsic
dual-energy condition characterizes all targets; every such positive
Radon measure is realized by an actual finite-energy metric, unique
modulo constants. No cardinality, toric, skeletal, support or
finite-generation restriction replaces the stated domain.

This edition combines the manuscript with the complete center,
local-product, common-collar and canonical-reference proofs
independently reconstructed by native review115. The earlier
manuscript's incomplete drafting status remains preserved in its
separate historical record. The only task-granted theorem is the exact
actual-ample P16.1 theory, used on explicitly ample projective auxiliary
carriers; the new cone, descent and local-operator applications are
proved here. Source-aware computational review is not a formal proof
certificate or a priority claim.

\section{1. Exact domain, conventions, and the ample
grant}\label{exact-domain-conventions-and-the-ample-grant}

Fix an algebraically closed field \(k\) of characteristic zero with its
trivial absolute value. Let \(X\) be a smooth integral proper
\(k\)-variety of dimension \(n>0\). Projectivity of \(X\) is not
assumed. By an actual rational line bundle we mean an element of
\(\operatorname{Pic}(X)\otimes\mathbb Q\), with all section calculations
made after taking specified integral multiples. A big actual rational
bundle \(L\) means that an integral multiple \(F=qL\) satisfies \[
\limsup_{m\longrightarrow\infty}m^{-n}h^0(X,mF)>0.
\] No section ring is assumed finitely generated. The construction
retains its choices; it does not identify theories from different
retained presentations.

The replacement datum will be \[
\mathcal M=(L,\pi:Y\longrightarrow X,A,D,\iota),\qquad
\iota:\pi^*L\simeq A+D,                                      \tag{1.1}
\] where \(Y\) is smooth integral projective, \(\pi\) is a proper
birational morphism, \(A\) is an actual ample rational line bundle, and
\(D\) is an effective rational Cartier divisor. The line-bundle
isomorphism is part of the datum. All choices are made before a target
measure is considered. Put \[V=(A^n)>0.\] This is the raw mass of the
genuine operator constructed below. We do not identify it with an
unrestricted minimal-singularity volume of \(L\), or with the possibly
different intersection number \((L^n)\).

\textbf{Main theorem.} Every \(X\) and specified big actual rational
\(L\) in Section 1 admits a datum \(\mathcal M\) as in (1.1). For each
retained datum, the entire dominated original-section cone
\(\mathcal P_{\mathcal M}\) of Section 5, its actual finite-energy space
\(\mathcal E^1_{\mathcal M}\) and the local nonpluripolar operator of
Sections 8 and 10 have normalized compactness, envelopes, decreasing
mixed and integrable weighted continuity, comparison and domination.
Their raw diagonal mass is \(V\). The intrinsic target condition of
Section 7 is finite exactly for the positive Radon measures of mass
\(V\) integrating every genuine relative finite-energy test. Each such
measure is the genuine operator of an actual
\(u\in\mathcal E^1_{\mathcal M}\), unique modulo constants. The
variational objective is attained, and its outer energy has the
two-sided derivative in every original \(C^0(X^{\rm an})\) direction.
The choices in \(\mathcal M\) are retained; compatibility between
distinct prescribed types is not part of this assertion.

\textbf{G0, used only on \((Y,A)\).} The task grants the genuine
finite-energy theory on geometrically unibranch integral projective
carriers with an actual ample rational bundle over the present fields.
In particular it grants the following on the smooth projective \(Y\):

\begin{itemize}
\tightlist
\item
  genuine \(A\)-psh potentials, finite fixed-bundle section
  approximation, divisorial determination, normalized compactness, and
  envelopes;
\item
  the genuine raw symmetric mixed intersection products for
  \(\mathcal E^1(Y,A)\), of mass \(V\), with locality, comparison,
  domination, and unbounded decreasing-net continuity;
\item
  integrability of finite-energy tests for mixed finite-energy products,
  and integrable weighted decreasing limits;
\item
  raw energy \(E_A\), its usc and monotone continuity properties, and
  its two-sided outer derivative in every direction in
  \(C^0(Y^{\rm an})\);
\item
  the intrinsic ample target condition, its equivalence to integrability
  of all actual \(\mathcal E^1(Y,A)\) potentials, actual variational
  attainment, and uniqueness modulo constants.
\end{itemize}

For a bounded genuine ample potential the granted raw energy has the
explicit formula \[E_A(\psi)=\frac1{n+1}\sum_{j=0}^n
\int_Y\psi\,(\omega_A+dd^c\psi)^j\wedge\omega_A^{n-j},\qquad
\omega_A=c_1(A)\text{ with its canonical zero metric}.\] For an upper
bounded unbounded genuine \(\psi\), it is the decreasing limit of these
energies for \(\max(\psi,-j)\); value \(-\infty\) is permitted. The
genuine finite-energy space consists exactly of potentials with finite
raw energy. These are the actual-ample definitions imported from G0,
before any transport to the original carrier. We use raw energy,
normalized by \[E_A(0)=0,\qquad E_A(\psi+c)=E_A(\psi)+Vc.\] The
probability calibration is \(e_A=V^{-1}E_A\). G0 is not a grant of
nonample varying-class closure, divisor division, nonprojective descent,
or a new local metric construction. Whenever its actual-ample
consequences are used below, the carrier is \(Y\) and the bundle is
\(A\), or another explicitly chosen actual ample bundle on \(Y\).

All analytic spaces are Berkovich analytifications. A Radon measure is a
regular Borel measure. Nets, compactness, and continuous test functions
are used without metrizability assumptions.

\section{2. Two approximation diagnostics and their
repairs}\label{two-approximation-diagnostics-and-their-repairs}

\subsection{2.1. The supplied counterexample is a proof, not an equation
obstruction}\label{the-supplied-counterexample-is-a-proof-not-an-equation-obstruction}

Let \(K\) be the disjoint union of an isolated point \(q\) and the
compact convergent sequence \(\{p_j:j\geq1\}\cup\{p_\infty\}\). Set
\(S=\mathbb N_{>0}\), \(T=\mathbb R\), and \(\rho(a)=1/a\). For
\(a,m\geq1\) and \(c\in\mathbb R\) let \[
f_{a,m,c}(q)=c,\quad f_{a,m,c}(p_a)=-\infty,\quad
f_{a,m,c}(x)=c-m\quad\text{at all remaining points}.
\] Each function is continuous into \(\mathbb R\cup\{-\infty\}\),
including at \(p_\infty\), and each family \(F_a\) is stable under
adding constants. For fixed \(a\), \(f_{a,m,0}\) decreases to the same
function \(g\): its value is zero at \(q\) and \(-\infty\) at every
other point. Thus \(g\in\widetilde F_{1/a}\) for every \(a\), and the
constant outer family \(g_a=g\) has parameters tending to zero.

Nevertheless \(\widetilde F_0\) is empty. In any decreasing net of
members of the \(F_a\), fix one term from \(F_a\). Every later term must
also be from \(F_a\), since its value at \(p_a\) must remain
\(-\infty\), whereas every member of every other family is finite there.
On that entire tail its parameter is \(1/a\), so the parameters cannot
converge to zero.

This disproves BJ v4 Lemma 4.6 in its printed generality. Its
strict-shift step fails at a common \(-\infty\) value; the corresponding
exponential difference is zero rather than strictly negative. The proof
does not disprove any specialized algebraic FS closure theorem, and
gives no counterexample to the desired equation theorem. We do not use
the lemma, or the nonample conclusions of BJ 4.15, 6.17, and 6.24
obtained through it.

\subsection{2.2. Safe finite fixed-bundle
flattening}\label{safe-finite-fixed-bundle-flattening}

Let \(B\) be an actual ample rational bundle on an integral projective
variety in the present field domain. Write \(F_B\) for finite,
real-valued, continuous actual FS section maxima \[
m^{-1}\max_j(\log\|s_j\|_0+\lambda_j),
\] where the displayed sections generate \(mB\) everywhere. Constants
and finite maxima preserve \(F_B\), after clearing degrees. Its zero
function is obtained by giving a generating family zero weights: at a
valuation-ring center every coefficient has norm at most one, and one
generator is a unit there.

Let \(\mathcal D_B\) consist of generically finite decreasing-net limits
from \(F_B\). Suppose \(\psi_i\downarrow\psi\) and
\(\psi_i=\inf_j f_{i,j}\) with \(f_{i,j}\in F_B\). Order triples
\((i,j,\epsilon)\) by increasing \(i\) and decreasing
\(f_{i,j}+\epsilon\), with \(\epsilon>0\). For two old triples, choose
\(i_3\) above both and \(\epsilon_3\) smaller than both old positive
errors. Then \[
\psi_{i_3}+\epsilon_3<
\min(f_{i_1,j_1}+\epsilon_1,f_{i_2,j_2}+\epsilon_2)
\] at every point. The right side is finite, so the inequality is still
strict where \(\psi_{i_3}=-\infty\). The exponentials of the inner
approximants minus the exponential of the finite right side decrease to
a strictly negative function. Their negative open sets cover the compact
analytic space; a finite subcover and directedness give one inner index
that works everywhere. Hence the triples are directed. The resulting net
\(f_{i,j}+\epsilon\) decreases to \(\psi\); at a pole use arbitrary
finite thresholds instead of subtracting infinities.

This proves the needed fixed-bundle nested-limit stability with
\textbf{finite inner approximants}. It gives no flattening theorem for
the common-pole original families after \(D\) is added. A single section
logarithm on \(B\) is in \(\mathcal D_B\): its floors are obtained by
adjoining a sufficiently divisible globally generating family of weight
\(-j\).

The fixed-bundle ample approximation granted by G0 therefore permits a
single finite \(F_A\) net for each genuine
\(\psi\in\operatorname{PSH}(Y,A)\). No countable extraction is required.

\subsection{2.3. Divided nets need not
decrease}\label{divided-nets-need-not-decrease}

Already on \(\mathbb P^1\), let \(\sigma\leq0\) be the logarithm of a
point-divisor section and consider \[
h_i=(1-i^{-2})\sigma,\qquad t_i=1-i^{-1}\quad(i\geq2).
\] The \(h_i\) decrease and \(h_i\leq t_i\sigma\), but their literal
quotients are \[h_i-t_i\sigma=(i^{-1}-i^{-2})\sigma.\] The positive
coefficient decreases for \(i\geq2\), so the quotients increase wherever
\(\sigma<0\). Their bundle classes tend to the ample remainder. Thus a
quotient-net monotonicity assertion cannot replace ample compactness.
Section 5 uses compactness instead.

\section{3. Existence of the retained algebraic
datum}\label{existence-of-the-retained-algebraic-datum}

\subsection{3.1. A smooth projective modification over every field in
the
domain}\label{a-smooth-projective-modification-over-every-field-in-the-domain}

Chow's lemma, in the form Stacks 0200 and 0201, gives an integral proper
birational refinement \(Z\to X\) with an immersion
\(Z\to\mathbb P^N_k\). Take the reduced closure of the common dense
open. Since \(Z\) is proper over \(k\), the immersion is closed, so
\(Z\) is projective and integral.

An integral finite-type scheme over any characteristic-zero field is
Noetherian, excellent, generically reduced, and a scheme over
\(\operatorname{Spec}\mathbb Q\). It need not be of finite type over
\(\mathbb Q\). The exact scope of Temkin's Theorem 1.2.1,
arXiv:0904.1592v2, applies to this \(Z\) and supplies a finite
regular-final blowup sequence. Integral blowup charts are subrings of
the common function field. Each blowup is projective by its degree-one
generated Rees algebra; composing the finite sequence retains
projectivity and birationality. The final \(Y\) is integral, regular,
projective, and of finite type over \(k\). Characteristic zero makes
\(k\) perfect, and Stacks 0B8X identifies regularity with smoothness in
this domain. Set \(\pi:Y\to X\) to be the composite.

Stacks 0AY8 applies with normal integral base \(X\), reduced integral
source \(Y\), and generic fiber \(\operatorname{Spec}k(X)\). It gives
\[\pi_*\mathcal O_Y=\mathcal O_X.                         \tag{3.1}\]
Locally trivialize any original bundle and apply (3.1) on that original
open. Gluing then gives, for each integral original bundle \(F\),
\[H^0(Y,\pi^*F)=H^0(X,F),                              \tag{3.2}\] with
the actual pullback map, not an abstract equality of dimensions.

\subsection{3.2. Big actual bundles exist on every original
carrier}\label{big-actual-bundles-exist-on-every-original-carrier}

Choose an effective very ample Cartier divisor \(H\) on \(Y\) containing
none of the finitely many exceptional prime divisors of \(\pi\). The
infinite field \(k\) permits a section avoiding this finite set of
linear conditions. Its pushforward \(B=\pi_*H\) is an effective Weil
divisor on \(X\). Smooth local rings are regular local UFDs (Stacks
0AG0), so \(B\) is an actual Cartier divisor. Birational equality away
from the exceptional locus gives \[\pi^*B=H+E,\] where \(E\) is
effective and exceptional: the pullback of an effective Cartier divisor
is effective, and \(H\) contains no exceptional prime. Multiplication by
\(s_E^m\) and (3.2) inject \[H^0(Y,mH)\hookrightarrow H^0(X,mB).\] The
left dimension grows with positive leading term of degree \(n\). Hence
\(\mathcal O_X(B)\) is an actual big line bundle. No ample bundle on the
original \(X\) has been introduced.

\subsection{3.3. Every specified big rational bundle has a polar
presentation}\label{every-specified-big-rational-bundle-has-a-polar-presentation}

Now fix any actual big rational \(L\) on \(X\) and clear denominators to
\(F=q\pi^*L\) on \(Y\). By (3.2), \(F\) has the same positive
section-growth limsup as \(qL\). Choose any actual very ample \(M\) on
\(Y\), and a nonzero hyperplane section \(T\in|M|\). We do not require
\(T\) to be reduced or smooth.

On the projective scheme \(T\), choose an ample line bundle \(N\)
sufficiently positive that \(N-F|_T\) is globally generated. A section
of that difference avoiding all the finitely many associated points is
regular. Such a section exists over the infinite field \(k\), since the
prohibited sections form a finite union of proper linear subspaces. Its
powers inject \(F^m|_T\) into \(N^m\). Serre vanishing and the coherent
Euler-polynomial support bound give
\[h^0(T,mF|_T)\leq h^0(T,mN)=O(m^{n-1}).\] The restriction sequence for
the effective Cartier \(T\) implies
\[h^0(Y,mF-M)\geq h^0(Y,mF)-O(m^{n-1}).\] Along arbitrarily large
integers realizing the positive \(m^n\) limsup, this is positive. Fix
one such \(m\) and one actual nonzero section with divisor \(G\geq0\).
Then \[\pi^*L=(qm)^{-1}M+(qm)^{-1}G.\] Retain \(A=(qm)^{-1}M\),
\(D=(qm)^{-1}G\), and this actual isomorphism. This proves (1.1) for
every specified big rational \(L\), using no Bertini statement or finite
generation.

\subsection{3.4. Whole bad fibers, rather than selected exceptional
points}\label{whole-bad-fibers-rather-than-selected-exceptional-points}

Let \(U_\pi\) be the maximal algebraic open over which \(\pi\) is an
isomorphism, and put \(B_\pi=X\setminus U_\pi\). Stacks 0BFP implies
that \(B_\pi\) has codimension at least two. More particularly, if there
is a quasi-finite point over a normal-base point, that base point
belongs to \(U_\pi\). Thus every irreducible component of a geometric
bad fiber has positive dimension: a zero-dimensional component would
contain a quasi-finite point. This argument applies after field
extension, since the smooth base remains normal and the birational
morphism remains proper.

On a reduced integral component \(F_0\) of such a geometric fiber,
\(\pi^*L\) is literally trivial, while \(A|_{F_0}\) is ample.
Consequently \(\mathcal O(mD)|_{F_0}\simeq\mathcal O(-mA)|_{F_0}\) for a
cleared positive integer \(m\). An antiample line bundle on a
positive-dimensional proper integral variety has no nonzero section.
Indeed a nonzero section and its powers would inject the unbounded
spaces \(H^0(F_0,rmA)\) into the finite space \(H^0(F_0,\mathcal O)\), a
contradiction. The canonical divisor section therefore vanishes on every
such component. We have proved \[\pi^{-1}(B_\pi)\subset |D|.\]

Let \[P=\pi(|D|)\subset X.\] This is a proper closed set: properness
makes the image closed, and the original generic point has its generic
singleton preimage outside \(D\). It contains \(B_\pi\). Outside
\(B_\pi\) every fiber is a singleton, and over \(B_\pi\) the entire
fiber lies in \(D\). It follows that \[\pi^{-1}(P)=|D|,
\qquad Y\setminus|D|\simeq X\setminus P.                 \tag{3.3}\] The
identities hold as closed supports, which is all that seminorms and
measures use. Analytification of closed supports and of the open
isomorphism gives \[p^{-1}(P^{\rm an})=|D|^{\rm an},\quad
p:Y^{\rm an}\setminus|D|^{\rm an}\simeq X^{\rm an}\setminus P^{\rm an}, \tag{3.4}\]
where \(p=\pi^{\rm an}\).

\section{4. Analytic centers, the original reference, and actual section
descent}\label{analytic-centers-the-original-reference-and-actual-section-descent}

\subsection{\texorpdfstring{4.1. Compactness and quotient topology
without projectivity of
\(X\)}{4.1. Compactness and quotient topology without projectivity of X}}\label{compactness-and-quotient-topology-without-projectivity-of-x}

For an analytic point \(x\), properness extends its associated morphism
\(\operatorname{Spec}\mathcal H(x)\to X\) uniquely to its valuation
ring. This gives the center \(c_X(x)\). Centers are compatible with
morphisms of proper varieties: compose the extended valuation-ring
morphism and use uniqueness. They are compatible with isometric
residue-field extensions for the same reason. These are actual centers,
not chosen valuations above the point.

Separatedness of algebraic \(X\) gives Hausdorffness of \(X^{\rm an}\)
by Berkovich 1993, 2.6.1 and 2.6.7. Compactness also has a direct center
proof. Choose a finite affine cover \(U_i=\operatorname{Spec}R_i\) of
\(X\). The set of seminorms on \(R_i\) bounded by one on every element
is a closed subset of \([0,1]^{R_i}\), since the multiplicative and
ultrametric identities are closed conditions. It is compact by
Tychonoff, without a cardinality restriction. Its image in
\(X^{\rm an}\) is exactly \(c_X^{-1}(U_i)\): boundedness supplies the
valuation-ring map, and conversely a center in \(U_i\) makes every
regular coefficient bounded by one. Those compact sets form a finite
cover of \(X^{\rm an}\).

The analytic morphism \(p\) is continuous and surjective. For
surjectivity one can use Berkovich 2.6.8, whose proof is by extension of
field valuations. In this particular proper setting, base change to
\(\mathcal H(x)\) and choose a closed point of the nonempty algebraic
fiber; its finite residue extension has the extended complete-field norm
and produces an analytic point above \(x\). Since \(Y^{\rm an}\) is
compact and \(X^{\rm an}\) Hausdorff, \(p\) is closed and is a quotient
map. Its compact fibers are the actual analytic fibers. Its restriction
to the open complements in (3.4) is the analytification of an actual
algebraic isomorphism. This proof does not use the omitted proof of
Berkovich 2.6.9.

\subsection{4.2. Canonical trivial metrics on the original
carrier}\label{canonical-trivial-metrics-on-the-original-carrier}

For an actual line bundle \(F\) on proper \(X\), define its canonical
trivial metric by making an algebraic frame at the valuation-ring center
have norm one. On \(c_X^{-1}(U_i)\), a frame \(e_i\) of \(F|_{U_i}\)
therefore has \(\|e_i\|_0=1\). On overlaps transition functions are
units at the center, hence have seminorm one. The metrics glue.

This metric is continuous. For any other algebraic frame \(e\) on an
open \(U\), its norm on each compact center set meeting \(U^{\rm an}\)
is the continuous seminorm of the corresponding transition function
\(e/e_i\). These sets give a finite relatively closed cover of
\(U^{\rm an}\). The closed-cover gluing criterion proves continuity on
\(U^{\rm an}\). Center compatibility proves the literal tensor and
pullback identities for canonical metrics, at every analytic point.
Clearing powers supplies the same statements for rational bundles. This
reconstructs the extension of the canonical-reference construction used
in BJ 1.24 to the present proper original carrier.

Write \[\sigma=\log\|s_D\|_0\leq0\] on \(Y^{\rm an}\), clearing the
rational divisor. It is continuous into \(\mathbb R\cup\{-\infty\}\) and
is \(-\infty\) exactly on \(|D|^{\rm an}\). For any genuine
\(\psi\in\operatorname{PSH}(Y,A)\), the function \(\sigma+\psi\) is usc,
bounded above, and constant on every fiber of \(p\). It is \(-\infty\)
on the whole preimage of \(P\), and all remaining fibers are singletons.
Let \[F(\psi)=u,\qquad p^*u=\sigma+\psi.                    \tag{4.1}\]
Closedness and surjectivity give \(\{u\geq a\}=p\{\sigma+\psi\geq a\}\)
closed for every finite \(a\), so \(u\) is usc. Its all-point value on
\(P\) is fixed, not inferred from divisorial density. At the generic
trivial point its value is \(\psi(v_{\rm triv})\), hence it is
generically finite. Put \(\tau=F(0)\); its pole set is exactly
\(P^{\rm an}\).

\subsection{4.3. A single genuine original-section
net}\label{a-single-genuine-original-section-net}

Choose a finite fixed-\(A\) net \(\psi_\alpha\downarrow\psi\) from
Section 2.2. At each stage clear all section degrees and bundle
denominators to \(r>0\). Multiply each section of \(rA\) by the
canonical section \(s_{rD}\). They are actual sections of \(r\pi^*L\).
By (3.2), each is the pullback of an actual section of \(rL\) on the
original \(X\). Canonical metric compatibility makes their logarithmic
norm formulas equal at \textbf{every} analytic point, with the same
constants. Let \(u_\alpha\) be their original maximum. Then
\[p^*u_\alpha=\sigma+\psi_\alpha,\qquad u_\alpha\downarrow F(\psi) \tag{4.2}\]
pointwise on \(X^{\rm an}\), including all poles. Neither the original
maxima nor \(\tau\) need be finite at their common base locus. The
flattening was performed before adding \(\sigma\); no strict inequality
at an original common pole was used. The reference \(\tau\) itself is a
finite original section maximum, since zero is a finite FS function on
\(A\).

These are actual singular metrics on \(L\): in an original local frame
\(e\), their raw weights are \(u_\alpha-\log\|e\|_0\), and their
decreasing limit is \(u-\log\|e\|_0\). They are not abstract energy
completions or an arbitrarily assigned lift.

\section{5. A direct repair of the entire fixed-divisor
cone}\label{a-direct-repair-of-the-entire-fixed-divisor-cone}

This section reconstructs the nonample foundations proposed in
continuation114. They are used only on the projective smooth \(Y\).
Define a numerical psh function by the original single-net definition:
it is a generically finite decreasing-net limit of actual
real-coefficient, generically finite FS maxima \(\phi_i\) in rational
bundles \(L_i\), with \(c_1(L_i)\to\theta\). No nested varying-class
closure is included as a premise.

\subsection{5.1. Rounding and embedding, directly in the defining
net}\label{rounding-and-embedding-directly-in-the-defining-net}

For positive rational \(\delta\), round the finitely many coefficients
of \(\phi_i\) upward to rational coefficients with error at most its
degree times \(\delta\). The resulting \(f_{i,\delta}\) is in the
\textbf{same} bundle and satisfies
\[\phi_i\leq f_{i,\delta}\leq\phi_i+\delta.\] Set
\(g_{i,\delta}=f_{i,\delta}+\delta\). Order pairs by increasing \(i\)
and decreasing \(g_{i,\delta}\). Given two pairs, take \(i_3\) above
both and \(2\delta_3\leq\min(\delta_1,\delta_2)\). Then
\[g_{i_3,\delta_3}\leq\phi_{i_3}+2\delta_3\leq
\min(\phi_{i_1}+\delta_1,\phi_{i_2}+\delta_2)\leq
\min(g_{i_1,\delta_1},g_{i_2,\delta_2}).\] These non-strict inequalities
hold at common poles as well. The pair set is directed; its projection
to the original index is final. Approximation at each finite value, or
below every real threshold at a pole, proves
\(g_{i,\delta}\downarrow\phi\) and preserves class convergence.

Choose an actual ample rational \(H\) with \(c_1(H)-\theta\) ample.
Eventually \(H-L_i\) is an actual ample rational bundle. Multiplication
by a zero-weight generating family of a divisible power of this
difference realizes \(\phi_i\) as an actual \(H\)-section maximum. If it
has base poles, adjoin a zero reference floor of weight \(-m\). The
product index \((i,m)\) gives a single decreasing net of finite
fixed-\(H\) maxima to \(\phi\). Thus \(\phi\) belongs to the genuine
actual-ample \(H\) cone. The ample G0 divisorial and usc-obstacle
principles now apply to it. A finite list of numerical functions can be
put in one common \(H\). This argument does not invoke nonample
regularization or varying-class closure.

\subsection{5.2. Homogeneous projection by a finite rational
polytope}\label{homogeneous-projection-by-a-finite-rational-polytope}

For \(\phi(v_{\rm triv})\geq0\) put
\[\widehat\phi(v)=\inf_{t>0}\frac{\phi(tv)}t.\] The two infima commute:
\[\inf_i\widehat\phi_i(v)=\inf_t\inf_i\phi_i(tv)/t=\widehat\phi(v).\]
For a rational finite FS expression \(m^{-1}\max_j(x_j+\lambda_j)\), set
\[\mathscr P=\{w_j\geq0:\ \sum w_j=1,\ \sum w_j\lambda_j\geq0\}.\] Then
\[\inf_{r\geq0}\max_j(x_j+r\lambda_j)
=\max_{w\in\operatorname{Vert}(\mathscr P)}\sum_jw_jx_j.       \tag{5.1}\]
Here is a direct one-dimensional proof. Every feasible weighted
expression bounds the left side from below. For finite \(x_j\), the
convex piecewise-linear function on the left attains its minimum at
zero, at a breakpoint with active slopes of opposite signs, or as a
limit at infinity. At a zero minimum an active nonnegative slope
supplies a unit-vector vertex. At a breakpoint, a zero weighted average
of two active opposite slopes supplies a vertex. In the final case the
maximal slope is zero and an active zero-slope unit vector supplies the
limiting value. These witnesses give equality. Vertices are the
indicated unit vectors and two-slope zero averages, hence rational. For
\(x_j=-\infty\), replace them by \(\max(x_j,-N)\) and commute the infima
over \(N\) and \(r\); the fixed finite maximum of vertex expressions
passes to its decreasing limit. Zero coefficients omit those coordinates
from a monomial.

Clear a common denominator \(b\) of the vertex weights. The products
\(\prod_j s_j^{bw_j}\) are actual sections of \(mbL\). Their homogeneous
logarithmic maximum divided by \(mb\) is \(m^{-1}\) times the right side
of (5.1), with all-point values including base poles. Hence
\(\widehat\phi_i\) are actual homogeneous rational FS functions in the
same bundles as \(\phi_i\), with generic value zero. They decrease to
\(\widehat\phi\), so the latter is genuine by the single defining net.
If \(h\) is already homogeneous, \(\widehat h=h\).

\subsection{5.3. Pointwise homogeneous decomposition by
convexity}\label{pointwise-homogeneous-decomposition-by-convexity}

Let \(C=\phi(v_{\rm triv})\) and \(h_\lambda=\widehat{\phi-\lambda}\)
for \(\lambda<C\). A tail of the defining net has nonnegative generic
values after subtracting \(\lambda\), so Section 5.2 applies. It gives
genuine homogeneous \(h_\lambda\), and \(h_\lambda+\lambda\leq\phi\) by
taking \(t=1\).

For a fixed actual point \(v\), the decreasing limit \(g(t)=\phi(tv)\)
is convex, decreasing, and bounded above by \(C\). It is either finite
on \((0,\infty)\) or identically \(-\infty\) there. To see the latter
assertion, a \(-\infty\) value at \(t_0\) and the common finite upper
bound at a smaller \(s\) force the limit to be \(-\infty\) at all
\(t\in(s,t_0)\) by convexity; monotonicity handles larger \(t\). Let
\(s\downarrow0\).

In the finite case choose a supporting line \(at+\lambda_0\) at \(t=1\).
Its slope is nonpositive, and \(\lambda_0\leq C\) by letting
\(t\downarrow0\) in \(at+\lambda_0\leq g(t)\leq C\). If \(\lambda_0<C\),
then \(h_{\lambda_0}(v)=a\), giving equality at \(v\). If
\(\lambda_0=C\), lower the intercept by \(\epsilon\):
\(h_{C-\epsilon}(v)\geq a\), giving the desired value within
\(\epsilon\). The identically infinite case is immediate. Therefore at
every actual point
\[\phi=\sup_{\lambda<C}(h_\lambda+\lambda).                \tag{5.2}\]
No usc replacement was made in this equality.

\subsection{5.4. Literal divisor division and ample
compactness}\label{literal-divisor-division-and-ample-compactness}

Let \(B=A+D\) be the retained actual bundle on \(Y\). Suppose \(h\) is
homogeneous numerical \(B\)-psh and \(h\leq\sigma\). Its homogeneous
actual FS approximants \(h_i\) from Section 5.2 decrease in bundles
\(B_i\) with classes tending to \(c_1(B)\). If \(D\ne0\), the level set
\(K=\{\sigma=-1\}\) is compact. For \(0<t<1\) rational, the open sets
\(\{h_i<-t\}\) increase to a cover of \(K\), so a tail satisfies that
inequality throughout \(K\). Homogeneity then gives \(h_i\leq t\sigma\)
at every function-field divisorial valuation with \(\sigma<0\), by
scaling to \(K\); when \(\sigma=0\), use \(h_i\leq0\). The common-ample
embedding and its usc-obstacle principle extend the inequality to all
points.

Use admissible pairs \((i,t)\) with \(i\) increasing and \(t\uparrow1\).
They are directed by first choosing a larger \(t\) and then a
sufficiently late \(i\). For each displayed section, the divisorial
inequality says its order along every prime divisor of \(D\) is at least
its degree times \(tD\). Clearing denominators and normality of \(Y\)
make division by \(s_{ptD}\) an actual global section of \(p(B_i-tD)\).
The resulting actual FS function \(r_{i,t}\) is the literal quotient on
the divisor complement and its actual extended value elsewhere; it has
generic value zero.

Refine by rational \(c>1\), tending to one, with \(cA-(B_i-tD)\) ample.
Openness of the ample cone gives admissible triples arbitrarily late and
with \(t,c\) arbitrarily close to one. Direct their indices by
increasing \(i,t\) and decreasing \(c\); a common upper triple is
obtained by fixing a smaller admissible \(c\) and taking the other
indices sufficiently late. Multiplying zero-weight generators of
\(A-(B_i-tD)/c\) realizes \(r_{i,t}/c\) in the actual \(A\) cone. If its
maxima have poles, the fixed-bundle floors from Section 5.1 give genuine
ample membership.

The normalized actual \(A\) cone is compact by G0. A convergent subnet
has a genuine quotient \(r\). At each divisorial point its values
converge explicitly to \(h-\sigma\), since \(h_i\), \(t\), and \(c\)
have the stated limits. Multiply a finite fixed-\(A\) net to \(r\) by
\(s_D\); this proves \(\sigma+r\) is genuine numerical \(B\)-psh. Both
it and \(h\) embed in one actual ample cone. Ample divisorial
determination now gives \(h=\sigma+r\) at every point. The quotient net
itself was not claimed decreasing. When \(D=0\), the same proof uses
\(h_i\) directly and omits \(K\) and division.

For an arbitrary numerical \(B\)-psh \(v\leq\sigma+C_0\), (5.2) gives
\(h_\lambda+\lambda\leq v\). Scaling any divisorial point and letting
the scale tend to infinity gives \(h_\lambda\leq\sigma\) there, hence
everywhere by the ample obstacle principle. Divide to get normalized
\(r_\lambda\). The family \(r_\lambda+\lambda\) is bounded above by
\(C=v(v_{\rm triv})\). Its ample usc envelope \(r\) is genuine. The
ample divisorial-nonnegligibility principle preserves its divisorial
values; thus (5.2) gives \(v=\sigma+r\) on divisorial points. Put both
sides in one ample cone and apply determination again. We have proved
the full all-point cone identity
\[\{v\in\operatorname{PSH}_{\rm numerical}(Y,B):v\leq\sigma+O(1)\}
=\sigma+\operatorname{PSH}(Y,A).                         \tag{5.3}\] The
quotient is unique, since \(\sigma\) is finite on function-field
divisorial points. Only actual-ample compactness, envelopes, and
determination were used; the nonample ingredients were reconstructed
from one defining FS net.

\subsection{5.5. Original section cone and its genuine
inverse}\label{original-section-cone-and-its-genuine-inverse}

Let \(\operatorname{PSH}_{\rm sec}(X,L)\) be the generically finite
single decreasing-net closure of actual original \(L\)-section maxima,
allowing their common base poles. Pulling its defining net to \(Y\)
gives a single defining numerical \(B\) net. Define the retained cone on
the original carrier by \[\mathcal P_{\mathcal M}
=\{u\in\operatorname{PSH}_{\rm sec}(X,L):u\leq\tau+C\text{ for some }C\in\mathbb R\}.\]
Section 4 proves
\(F(\operatorname{PSH}(Y,A))\subset\mathcal P_{\mathcal M}\). Conversely
the actual pullback of \(u\in\mathcal P_{\mathcal M}\) satisfies
\(p^*u\leq\sigma+C\), so (5.3) supplies its unique genuine quotient
\(Q(u)\). Surjectivity of \(p\) gives \(u=F(Q(u))\) at every point. Thus
\[F:\operatorname{PSH}(Y,A)\simeq\mathcal P_{\mathcal M},\qquad Q=F^{-1}. \tag{5.4}\]
This is genuine section descent, not an arbitrarily chosen lift. It
reflects order: compare on divisorial points, cancel their finite
\(\sigma\), and use the ample usc-obstacle principle. It commutes with
constants and finite maxima. Its divisorial topology is the original
pointwise topology on the common function field, translated by the fixed
finite values of \(\sigma\). Hence (5.4) is a homeomorphism, preserving
suprema at the generic trivial point. The normalized original cone is
compact; decreasing original limits have a genuine quotient and an
actual original section net by Section 4.

The numerical zero-class repair in continuation114 is also valid within
these foundations. Given a single defining FS net with classes tending
to zero, generic values \(C_i\downarrow C\), refine by
\(\epsilon\downarrow0\) with \(\epsilon A-L_i\) ample. The functions
\((\phi_i-C_i)/\epsilon\) belong to the normalized actual \(A\) cone. At
every divisorial point \(v\), compactness and continuous evaluation give
a common finite lower bound \(M_v\). Thus
\[\epsilon M_v\leq\phi_i(v)-C_i\leq0.\] The limit equals \(C\)
divisorially and, after one ample embedding, everywhere. This
establishes a constants-only zero class without applying varying-class
closure to \(m\phi\). It is not used to extend the present task to
zero-volume nonample classes.

\section{6. Original finite energy, fiber maxima, and a finite-energy
pole}\label{original-finite-energy-fiber-maxima-and-a-finite-energy-pole}

Define, intrinsically for the retained genuine cone,
\[E_{\mathcal M}(u)=E_A(Q(u)),\qquad
\mathcal E^1_{\mathcal M}=\{u\in\mathcal P_{\mathcal M}:E_{\mathcal M}(u)>-\infty\}. \tag{6.1}\]
This energy is attached to actual original metrics, with its genuine
quotient uniquely recovered in (5.4). It is concave, monotone, weakly
usc, and decreasing-net continuous by G0 and the order homeomorphism.
Its constant law is \(E_{\mathcal M}(u+c)=E_{\mathcal M}(u)+Vc\).

For \(u=F(\psi)\) set
\[q_u(x)=\max_{y\in p^{-1}(x)}\psi(y).                    \tag{6.2}\]
The maximum is attained on every compact fiber, including value
\(-\infty\). The sets \(\{q_u\geq a\}=p\{\psi\geq a\}\) are closed.
Hence \(q_u\) is an actual upper bounded usc function on the original
analytic carrier. On its open complement of \(P\) it equals \(u-\tau\).
On \(P\), formula (6.2) gives a uniquely specified value from the
retained genuine quotient; it is never an expression
\(-\infty-(-\infty)\). If \(\psi_i\downarrow\psi\), then
\(q_{F(\psi_i)}\downarrow q_{F(\psi)}\): on a fixed fiber the closed
sets \(\{\psi_i\geq a\}\) decrease, and compactness excludes a
persistent level exceeding the limiting maximum. Constants and order are
respected. Additivity on bad fibers is not asserted or needed.

There is an actual \(w\in\mathcal E^1(Y,A)\) equal to \(-\infty\) on all
of \(D\). Choose \(r\) sufficiently large that \(rA-D\) has an effective
nonzero section. Multiplying by the divisor section gives a section of
\(rA\) vanishing along \(D\). Its normalized logarithm \(a\leq0\) is
genuine \(A\)-psh by fixed-bundle floors. Set
\[\alpha=\frac1{n+2},\qquad w=1-(1-a)^\alpha.             \tag{6.3}\]
The convex increasing function in (6.3) has slopes in \([0,1]\), value
zero at zero, and value \(-\infty\) at \(-\infty\). Its affine supports
applied to \(a\) are genuine \(A\) potentials by convexity with zero.
Their bounded-above ample envelope is exactly the usc composition \(w\),
proving genuine membership without nonample composition closure.

Here is the finite-energy estimate, including its necessary exponent. In
probability calibration let
\[\operatorname{Cap}(E)=\sup_{-1\leq b\leq0,\ b\in\operatorname{PSH}(A)}\operatorname{MA}_{A,\rm prob}(b)(E).\]
For every normalized genuine \(a\),
\(\operatorname{Cap}\{a\leq-s\}\leq n/s\). To justify this without
unbounded integration assumptions, apply signed integration by parts to
the bounded floors of \(a\). Telescope the difference between the
measure of \(b\) and that of zero. Each of the \(n\) terms is an
integral of \(b\in[-1,0]\) against the difference of two positive mixed
products of the same probability mass, and has absolute value at most
one. The zero-potential measure is the generic trivial Dirac measure,
where every normalized floor vanishes. Monotone convergence in the fixed
measure of \(b\) gives
\(\int(-a)\operatorname{MA}_{\rm prob}(b)\leq n\). Chebyshev and the
supremum over \(b\) prove the capacity bound.

For any normalized genuine \(z\), write \(z_s=\max(z,-s)\),
\(\mu_s=\operatorname{MA}_{\rm prob}(z_s)\). Energy concavity gives
\(J(z_s)\leq\int(-z_s)d\mu_s\). Locality and equal masses show, for
\(0<t\leq s\), that \(\mu_s\{z\leq-t\}=\mu_t\{z\leq-t\}\). For
\(t\geq1\), positivity of the mixed expansion gives
\(\operatorname{MA}_{\rm prob}(z_t/t)\geq t^{-n}\mu_t\). Layer cake
therefore gives
\[J(z_s)\leq1+\int_1^\infty t^n\operatorname{Cap}\{z\leq-t\}\,dt. \tag{6.4}\]
The factor \(t^n\) is retained; it is missing in a printed intermediate
line in BJ's proof. In (6.3),
\(\{w\leq-t\}=\{a\leq1-(1+t)^{1/\alpha}\}\), so its capacity is
\(O(t^{-1/\alpha})\). The integrand in (6.4) is \(O(t^{-2})\). The
energies of the floors are bounded below uniformly, and decreasing
energy continuity proves \(w\in\mathcal E^1(Y,A)\). We make no claim for
the larger printed exponent interval in BJ 7.48.

Since \(p^{-1}(P)=D\), (6.2) gives \(q_{F(w)}=-\infty\) on all of \(P\).
Every mixed finite-energy product on \(Y\) integrates \(w\) by G0, hence
gives \(D\) zero mass. This establishes the exceptional-set conclusion
used below before any target is chosen.

\section{7. Radon inversion and the original intrinsic target
class}\label{radon-inversion-and-the-original-intrinsic-target-class}

We first prove the relevant measure theorem. Suppose \(p:S\to T\) is a
closed continuous surjection of compact Hausdorff spaces, \(C\subset T\)
is closed, \(p^{-1}(C)=F\), and \(p:S\setminus F\simeq T\setminus C\)
has inverse \(j\). Every Radon \(\mu\) on \(T\) with \(\mu(C)=0\) has a
unique Radon inverse \(\nu\) on \(S\) satisfying \(\nu(F)=0\) and
\(p_*\nu=\mu\).

For \(g\in C(S)\) define
\[\Lambda(g)=\int_{T\setminus C}g(j(x))\,d\mu(x).\] This is a positive
bounded functional, and Riesz gives a Radon \(\nu\). To prove zero mass
on \(F\), choose compact \(K\subset T\setminus C\) with
\(\mu(T\setminus K)\) arbitrarily small. A continuous \(\rho\in[0,1]\)
equal to one on \(C\) and zero on \(K\) satisfies
\(\nu(F)\leq\Lambda(\rho\circ p)=\int\rho d\mu\), which is arbitrarily
small. Continuous tests pulled from \(T\) prove pushforward. Uniqueness
follows from the isomorphism off \(F\), zero mass on \(F\), and
regularity; there is no measurable-selection or arbitrary-fiber choice.
Integrals of upper bounded usc functions agree, including \(-\infty\),
by their continuous decreasing majorants or ordinary Borel transport.

Apply this to \(S=Y^{\rm an}\), \(T=X^{\rm an}\), \(C=P^{\rm an}\), and
\(F=D^{\rm an}\). For a positive Radon \(\mu\) of mass \(V\), define the
original retained dual energy
\[D_{\mathcal M}(\mu)=\sup_{u\in\mathcal E^1_{\mathcal M}}
\left(E_{\mathcal M}(u)-\int q_u\,d\mu\right).             \tag{7.1}\]
The quotient integral is relative to the actual reference; absolute
integrability of \(u\) or \(\tau\) is not asserted. If (7.1) is finite,
every \(q_u\) is integrable: otherwise its individual expression would
be \(+\infty\). In particular the explicit pole \(q_{F(w)}\) forces
\(\mu(P)=0\). The unique Radon inverse \(\nu\) exists and
\[\int q_u\,d\mu=\int Q(u)\,d\nu,\qquad
D_{\mathcal M}(\mu)=E_A^\vee(\nu).                       \tag{7.2}\] G0
proves that (7.1) is finite if and only if \textbf{every} \(q_u\),
\(u\in\mathcal E^1_{\mathcal M}\), is \(\mu\)-integrable. In the reverse
implication, the pole again excludes \(P\) first, then G0 is applied
through (7.2). Thus this target condition is intrinsic to the retained
original metric cone and its reference. It specifies the exact mass and
energy and imposes no toric, skeletal, or countability restriction. Zero
mass on \(P\) is a proved consequence, rather than a silent support
assumption.

For completeness, zero-exceptional Radon inversion is weakly continuous
at zero-exceptional limits. If \(\mu_i\to\mu\) weakly, all masses equal
\(V\), and all \(\mu_i(P)=\mu(P)=0\), let \(\nu_i,\nu\) be their
inverses. For \(g\in C(Y)\) choose a bump \(\rho=1\) near \(P\) with
\(\int\rho\,d\mu\) small. The function \((1-\rho\circ p)g\) descends off
\(P\) and extends continuously by zero across \(P\). Its integrals
converge. The remaining integrals have absolute values at most
\(\|g\|_\infty\int\rho\,d\mu_i\), which are eventually small. This
proves weak convergence of the inverses for arbitrary nets.

\section{8. Pretarget mixed products from original section
metrics}\label{pretarget-mixed-products-from-original-section-metrics}

\subsection{8.1. The candidate pairing identified by the local
construction}\label{the-candidate-pairing-identified-by-the-local-construction}

For \(u_i=F(\psi_i)\in\mathcal E^1_{\mathcal M}\) define
\[\operatorname{MA}_{\mathcal M}(u_1,\ldots,u_n)
=p_*\bigwedge_{i=1}^n(c_1(A)+dd^c\psi_i).                \tag{8.1}\]
This pairing is constructed for the entire cone before any target is
chosen. It is a positive symmetric Radon measure of raw mass \(V\),
giving \(P\) zero mass by Section 6. Section 10 constructs the intrinsic
local products of the actual original metric and proves that their
nonpluripolar reference-floor extension equals this pairing. Thus it is
the genuine original-carrier operator; the equality with the ample
measure is a conclusion of that local construction.

G0 and compact pushforward give unbounded decreasing-net weak
convergence in every entry, with genuine finite-energy limits.
Integrable weighted limits also transport. The defining identity has the
sign \[\int\psi_0\,\bigwedge_{i=1}^n(c_1(A)+dd^c\psi_i)
=(c_1(A),\psi_0)\prod_{i=1}^n(c_1(A),\psi_i)
-(c_1(A),0)\prod_{i=1}^n(c_1(A),\psi_i).                 \tag{8.2}\] It
is this difference, not its reverse, that is passed to the limit. On the
original carrier, (8.2) gives
\(q_{u_0}\operatorname{MA}_{\mathcal M}=p_*(\psi_0 M_A)\). For
integrable decreasing tests and decreasing finite-energy mixed factors,
G0 gives the weighted weak limit. A difference of two such tests is
treated term by term; the difference is not falsely called decreasing.
Ordinary weak pushforward is legitimate on the compact spaces
established in Section 4.

\subsection{8.2. Literal local-frame signs at continuous section
stages}\label{literal-local-frame-signs-at-continuous-section-stages}

Let \(W\subset X^{\rm an}\setminus P^{\rm an}\) be an ordinary open
frame neighborhood with compact closure still inside a common frame
region. Pull it to \(Y\setminus D\) using the actual isomorphism (3.4).
Choose compatible frames \(e_A,e_D\) under \(\iota\), so the pulled
original frame is \(e_L=e_Ae_D\). Write \(s_D=f e_D\), where \(f\) is an
algebraic, hence analytic, unit on this region. Let
\(h_A=-\log\|e_A\|_0\), \(h_D=-\log\|e_D\|_0\). Then
\[\sigma=\log|f|-h_D,\qquad
u+h_L=\psi+h_A+\log|f|.                                \tag{8.3}\] The
unit logarithm has a \textbf{plus} sign. On a compact frame neighborhood
its value and the reference weights are finite and bounded.

For a finite \(F_A\) section maximum \(\psi_\alpha\), (4.2) gives the
actual original section maximum, with raw weight
\(\psi_\alpha+h_A+\log|f|\). Its finite-stage local metric intersection
product is the usual product for that actual section metric. BE 7.14
constructs smooth semipositive regularizations of the finite \(A\)
metric by regularized maxima. BE 8.4 gives the uniform-to-weak
continuous product and its identification with the genuine finite
intersection product. CLD 18.1.2 gives the zero Hessian of \(\log|f|\),
and CLD 8.8.5 gives its invariance at continuous local stages.
Restriction through the analytic isomorphism therefore identifies the
finite \textbf{actual original local} product with the finite ample
product on \(p^{-1}(W)\). This is a pretarget metric identity, not an
identity imposed by a desired measure.

We use those primary statements only in their continuous domains. They
do not establish a bounded discontinuous local product merely by taking
uniform limits.

\subsection{8.3. A precise decreasing local section
category}\label{a-precise-decreasing-local-section-category}

Here is the part that can be proved directly from G0. On the preceding
\(W\), consider nets of raw weights
\[\lambda_{i,\alpha}=\psi_{i,\alpha}+h_A+\log|f|,\] where
\(\psi_{i,\alpha}\) are finite actual \(F_A\) maxima, decreasing on all
of \(Y^{\rm an}\) to a genuine bounded \(A\)-psh endpoint \(\psi_i\).
The common directed index is obtained by a finite product of index sets.
Call these \emph{divided global-section nets on \(W\)}. They are actual
original-section local weights by (4.2).

For \(g\in C_c(W)\), pull \(g\) through the open isomorphism and extend
it by zero on \(Y\). G0's bounded decreasing-net continuity proves
convergence of its integrals against the finite local intersection
products in Section 8.2. The resulting local Radon measure is
\[M_A(\psi_1,\ldots,\psi_n)|_{p^{-1}(W)},                \tag{8.4}\]
read on \(W\) through the isomorphism. This conclusion applies even when
every \(\psi_i\) is bounded but discontinuous. It uses monotone ample
product convergence, not uniform convergence to the endpoint.

The limit is independent of the particular divided global-section nets.
If two bounded global endpoint tuples have equal raw weights on an
original open \(W\), their fixed reference and unit terms in (8.3)
cancel, leaving equal ample scalar potentials there. G0's mixed locality
makes (8.4) equal on \(W\). This also proves independence from changes
of the global endpoints outside \(W\), within this category. Frame
changes are handled at \textbf{each} finite stage: changing the original
frame subtracts the logarithm of its analytic unit from the raw weight,
without changing the finite local product. The same transformed
approximants pass through the decreasing construction, so their limits
agree as well. The fixed unit term in (8.3) has thus passed through the
bounded construction; its zero-Hessian property was not simply asserted
at the discontinuous endpoint.

One can also see the requisite ample locality directly from the granted
strict-winning identity. For equal potentials \(a,b\) on \(W\), apply
strict-winning locality to \(a+\epsilon,b\) there, then let
\(\epsilon\downarrow0\) in the genuine decreasing products. Exchange
\(a,b\) and polarize. All entries remain in the actual ample G0 domain.
Weighted versions use (8.2) with two separate decreasing terms.

This proves a genuine decreasing local limit for the full
type-preserving original-section category, with all-point endpoints and
germ independence. Section 10 proves its equality with independently
supplied admissible local contact nets using one common actual ample
collar.

\subsection{8.4. Relative floors and the unbounded nonpluripolar
product}\label{relative-floors-and-the-unbounded-nonpluripolar-product}

For \(u_i=F(\psi_i)\in\mathcal E^1_{\mathcal M}\) put
\[u_{i,j}=\max(u_i,\tau-j)=F(\max(\psi_i,-j)).             \tag{8.5}\]
These are genuine bounded-quotient metrics. On compact original frame
neighborhoods off \(P\), their raw weights are bounded, possibly
discontinuous, and belong to Section 8.3. Absolute original floors
\(\max(u_i,-j)\) are not substituted for (8.5).

For the mixed product let \(G_j=\bigcap_i\{\psi_i>-j\}\). G0's mixed
strict-winning locality, obtained by polarization, gives
\[\mathbf1_{G_j}M_A(\max(\psi_1,-j),\ldots,\max(\psi_n,-j))
=\mathbf1_{G_j}M_A(\psi_1,\ldots,\psi_n).                \tag{8.6}\]
Each \(\psi_i\) is integrable for the latter measure, so its \(-\infty\)
set has zero mass. The measure also gives \(D\) zero mass. The
restrictions of the left side to \(Y\setminus D\) therefore increase to
the full genuine ample mixed measure. Their original restrictions give a
Radon measure of mass exactly \(V\), equal to (8.1), and zero extension
across \(P\) is unique for a measure required to charge that pole set by
zero.

The restricted limit is independent of floor levels tending to infinity.
More generally, a bounded genuine \(A\)-psh reference \(b\) can replace
zero in \(\max(\psi_i,b-j)\); the same strict-winning and integrability
argument gives the same restriction. Weighted limits for integrable
genuine quotient tests follow from (8.2) and G0. Fixed frame-unit
transformations continue to agree at the bounded-floor stages and hence
in (8.6).

These assertions prove the relative nonpluripolar extension in the
original-section category, including mass and weighted limits. The
independently constructed local product in Section 10 agrees with these
restrictions and floors, completing the intrinsic operator
identification before any target is chosen.

\section{9. Original envelopes, all continuous variations, and
equations}\label{original-envelopes-all-continuous-variations-and-equations}

Section 10 identifies the pretarget mixed pairing with the intrinsic
local original-carrier operator. The following deductions therefore give
the equation theorem for that genuine operator.

\subsection{9.1. Genuine minorants and all original continuous majorants
agree}\label{genuine-minorants-and-all-original-continuous-majorants-agree}

For an upper bounded original obstacle \(h\), let
\[\mathscr O_{\mathcal M}(h)=\sup\{E_{\mathcal M}(v):v\in\mathcal P_{\mathcal M},\ v\leq h\}.\]
Take \(u=F(\psi)\in\mathcal E^1_{\mathcal M}\) and
\(f\in C^0(X^{\rm an})\). Cancel finite \(\sigma\) at function-field
divisorial points and apply the actual ample usc-obstacle principle.
This proves the exact equivalence
\[F(\eta)\leq u+tf\text{ on }X^{\rm an}\quad\Longleftrightarrow\quad
\eta\leq w_t:=\psi+t p^*f\text{ on }Y^{\rm an}.           \tag{9.1}\]
The pole fibers give no lost condition, because the divisorial
inequalities determine the genuine ample potentials everywhere. The
ample envelope \(P_A(w_t)\) lies between \(\psi\pm |t|\|f\|_\infty\),
hence is actual finite energy. Thus
\[\mathscr O_{\mathcal M}(u+tf)=E_A(P_A(w_t))=E_A^\downarrow(w_t). \tag{9.2}\]

There is also an outer definition using \textbf{all original continuous
relative majorants}: \[\mathscr O^{\rm out}_{\mathcal M}(u+tf)
=\inf_{b\in C^0(X^{\rm an}),\ q_u+tf\leq b}
E_{\mathcal M}(F(P_A(p^*b))).                            \tag{9.3}\] For
any \(b\in C^0(X)\), (9.1) and the fiber maximum show
\[u+tf\leq\tau+b\quad\Longleftrightarrow\quad q_u+tf\leq b.\] Compact
Hausdorff normality supplies continuous majorants of any upper bounded
usc function whose pointwise infimum is that function; finite minima
make the family a decreasing net. Let \(b_i\downarrow q_u+tf\) be this
net of all original majorants, and put \(\theta_i=P_A(p^*b_i)\). These
are genuine and decrease, bounded below by \(P_A(w_t)\) and above by a
fixed continuous bound. G0 gives a genuine finite-energy limit
\(\theta\).

Every function-field divisorial point has a unique analytic birational
lift: apply the proper valuative criterion to its common function-field
valuation ring. Its support is the original generic point, so \(p\) has
a singleton analytic fiber there and \(q_u(p(v))=\psi(v)\). It follows
that \(\theta(v)\leq w_t(v)\) for every such divisorial point. Ample usc
determination extends this inequality to all \(Y\). Conversely
\(\theta_i\geq P_A(w_t)\), since \(p^*b_i\geq w_t\). Hence
\(\theta=P_A(w_t)\). Decreasing energy continuity gives
\[\mathscr O^{\rm out}_{\mathcal M}(u+tf)=\mathscr O_{\mathcal M}(u+tf). \tag{9.4}\]
Only original continuous majorants were used in (9.3); no contraction to
a projective quotient restricts the variations. Continuous relative
obstacles \(\tau+b\) have actual envelopes \(F(P_A(p^*b))\), continuous
relative to the reference on the divisor complement, and with the
prescribed all-point poles.

\subsection{9.2. Both derivative signs in every original
direction}\label{both-derivative-signs-in-every-original-direction}

Applying G0 to \(w_t=\psi+t p^*f\) and (9.2)--(9.4) gives
\[\left.\frac d{dt}\right|_{t=0}\mathscr O_{\mathcal M}(u+tf)
=\int_Y p^*f\,M_A(\psi)=\int_X f\,\operatorname{MA}_{\mathcal M}(u). \tag{9.5}\]
This is a two-sided derivative for every \(f\in C^0(X^{\rm an})\), not
just directions descending from a projective contraction or PL tests.
Its factor is raw mass \(V\). For clarity, when extending a PL
derivative by uniform approximation with error \(\epsilon\), the
obstacle bounds are \(\pm|t|\epsilon\) and the quotient errors are
\(\pm V\epsilon\). Replacing \(f\) by \(-f\) gives the second sign; no
\(t^2\epsilon\) bound is used.

\subsection{9.3. An actual attained maximizer and the
equation}\label{an-actual-attained-maximizer-and-the-equation}

Let \(\mu\) be a positive Radon measure of mass \(V\) satisfying (7.1).
Its unique Radon inverse \(\nu\) has finite ample dual energy by (7.2).
G0 gives an actually attained genuine \(\psi\in\mathcal E^1(Y,A)\)
maximizing \(E_A(\eta)-\int\eta d\nu\). Hence \(u=F(\psi)\) is an actual
original metric attaining (7.1), not a completion or an approximate
object.

The mechanism can be expressed on the original carrier. A normalized
maximizing family has bounded ample \(J\), hence bounded original \(J\)
under (5.4). Original normalized compactness gives a subnet to an actual
\(u\); energy usc and the ample bounded-\(J\) target \(L^1\) continuity
through (7.2) give attainment. Nonmetrizable compactness is used as
subnet compactness. No countability or countable field containing an
entire psh family is required.

For every original continuous \(f\) and real \(t\), the genuine minorant
\(v_t=F(P_A(w_t))\) satisfies \(q_{v_t}\leq q_u+tf\), by order and
(9.1). Integrability and the maximal objective therefore give
\[\mathscr O_{\mathcal M}(u+tf)-\int(q_u+tf)d\mu
\leq E_{\mathcal M}(u)-\int q_ud\mu.\] Use (9.5), first with \(t>0\) and
then with \(t<0\), reversing the inequality when dividing by a negative
number. We obtain equality of the two finite Radon measures on every
continuous original test:
\[\operatorname{MA}_{\mathcal M}(u)=\mu.                       \tag{9.6}\]
By the local identification proved in Section 10, this is the genuine
original-carrier finite-energy Monge--Ampere equation. The attained
solution is an actual original-section metric.

For two solutions of (9.6), their ample measures have zero mass on \(D\)
and identical pushforwards. Radon inversion proves equality of those
ample measures; G0 uniqueness makes their genuine quotients differ by
one constant. Thus their actual original metrics differ by the same
constant at every point, including their prescribed poles. The converse,
that such a solution has finite dual energy and computes the dual
supremum, follows from the ample energy concavity statement through
(7.2).

\subsection{9.4. Comparison and domination for the genuine
operator}\label{comparison-and-domination-for-the-genuine-operator}

For \(u,v\in\mathcal E^1_{\mathcal M}\), the strict sets outside \(P\)
correspond exactly to the strict quotient sets outside \(D\), and both
product measures charge \(P\) by zero. Ample comparison therefore gives
\[\int_{\{u<v\}}\operatorname{MA}_{\mathcal M}(v)
\leq\int_{\{u<v\}}\operatorname{MA}_{\mathcal M}(u).            \tag{9.7}\]
Strict-winning locality and mixed locality transport in the same way. If
a genuine \(u\in\mathcal P_{\mathcal M}\) is at most a finite-energy
\(v\) almost everywhere for \(\operatorname{MA}(v)\), Radon inversion
transfers that inequality off \(D\) to the ample measure. It gives \(D\)
zero mass, so ample domination applies and yields \(Q(u)\leq Q(v)\)
everywhere. Then \(u\leq v\) everywhere. None of these deductions
assumes projectivity of the original carrier or a prescribed skeletal
support.

\section{10. Intrinsic local products through a common actual ample
collar}\label{intrinsic-local-products-through-a-common-actual-ample-collar}

We now prove the category comparison needed to identify the pairing of
Section 8 with the intrinsic original metric operator. Admissible local
inputs are decreasing finite affine FS/contact weights on a common
ordinary neighborhood, with the actual metric cocycles and bounded
endpoint. The associated entire original cone is recovered by its
reference floors; no density theorem for arbitrary unrelated classical
currents is asserted. Appendix D gives the finite
contact-to-genuine-ample reconstruction used in the proof. Appendix C
proves continuous canonical original references and signed frame
compatibility.

Retain the datum and genuine cone conclusions of Section 5 and Appendix
B. The field is any algebraically closed trivially valued
characteristic-zero k. In this section all local approximants are actual
finite rational affine FS weights on one ordinary neighborhood, with
bounded decreasing endpoints. Real constants can be handled by the
direct upward-rounding net explained below. No arbitrary classical CLD
psh/current density theorem is assumed. The complete associated original
cone is not reduced: its bounded reference floors and actual local
section approximants have exactly these weights, and all E1 endpoints
are recovered from their floors.

\subsection{10.1. The only local bridge, applied on an actual ample
projective
carrier}\label{the-only-local-bridge-applied-on-an-actual-ample-projective-carrier}

The exact G0 mathematical body, Section2, states the compact contact
theorem and its continuous contact metric bridge in its projective,
actual-ample domain. The continuous bridge concerns an ALREADY GLUED
continuous rational metric on an actual ample rational H, whose local
weights have finite unit-log affine supports below the weight, equal at
the point with the SAME additive constant. It makes that metric
genuinely H-psh. This is the precise input used below, ONLY on the
smooth projective Y and an actual ample H. It is not applied to the
original nonprojective X or its nonample L.

Appendix D includes the complete pertinent reconstruction of this
bridge, using actual algebraic PL extensions, represented models and
same-bundle Gauss descent. Its new application here is established from
that proof, rather than granted by a historical label. It first
constructs exact original algebraic PL extensions using clipped section
logarithms, actual principal localization and finite rational polyhedral
subdivision. Flag ideals give an algebraic test configuration on the
SAME actual H. Over k((t)) the pulled unit-log supports give the curve
maximum test. The GR and GM model criteria test whole proper residual
curves and imply nefness of that supplied algebraic model. BE7.8 is used
only over this NONTRIVIAL extension. Gauss evaluation of its pure FS
sections gives an EXACT finite coefficient cutoff in original H
sections: all coefficient norms are at most1 and the continuous metric
has a compact lower bound. Uniform FS approximants with summable shifts
decrease in that fixed H.

This lists the actual route and its primary background; it is not a new
proof of every recursively cited theorem. Those primary bodies are in
the supplied G0 source dictionary. Every new application must check
their exact hypotheses. No general varying-class/pole flattening lemma,
arbitrary formal-line-bundle algebraization, Laurent-cover positivity
criterion, or unrestricted current converse is used by this bridge.

\subsection{10.2. Exact common geometric neighborhood and
reference}\label{exact-common-geometric-neighborhood-and-reference}

Fix a point of \(V=X^{\rm an}\setminus P^{\rm an}\) and identify its
neighborhood with an ordinary open of
\(Y^{\rm an}\setminus D^{\rm an}\). Shrink to a principal algebraic
chart where the original L frame, the A frame and D section factor are
defined and the D factor is a unit. Basic seminorm neighborhoods permit
finitely many strict rational upper/lower bounds on regular functions;
required nonvanishing functions are inverted in this same chart. Append
a complete very ample section system so that its ratios include a
generating set of the original affine chart ring. These generator bounds
make the closed seminorm domain compact. They do not require strict
affinoidness over the trivial field or a countable field.

The explicit construction in Appendix D.4 clears the principal
denominators and builds ONE actual FS local reference weight v on an
actual rational ample H, with rational epsilon\textgreater0 and

\begin{verbatim}
Q={v<=0} compactly contained in the given ordinary neighborhood,
P0={v<=-epsilon}, and the point in interior(P0).
\end{verbatim}

The local raw expression, with tensor denominator q, is

\begin{verbatim}
v=(1/q) max(q lambda, log|F_j/a^m|-C_j).
\end{verbatim}

The thresholds are C\_j, without an erroneous extra factor q. A negative
constant branch and the appended section system give a finite weight.
Clearing its denominator produces actual H sections. Write v=h0+p, where
h0 is the canonical H raw weight in that frame and p is the
corresponding GLOBAL actual finite FS potential. It is the same p and
the same H for every approximation to be compared. Local h0 is finite
continuous on Q, and all cocycles are those of that actual H.

For compact tests it is enough to make this construction at each point
of their support and take a finite cover. There is no claim that one
principal chart contains an arbitrary large compact set.

\subsection{10.3. One fixed collar for finitely many nets, including
independent
nets}\label{one-fixed-collar-for-finitely-many-nets-including-independent-nets}

Let w\_i be a finite tuple of bounded raw local weights on a common
ordinary neighborhood of Q. Each is the pointwise decreasing limit of
finite rational AFS/contact weights w\_(i,alpha). Include BOTH nets for
any one endpoint to be compared, as well as a bounded test weight if
needed. Choose rational M\textgreater0 so -M\textless=w\_i\textless=M on
Q. Choose rational R\textgreater0 so R epsilon\textgreater2M+2, and set

\begin{verbatim}
h=R v+M+1.
\end{verbatim}

Then h\textless-M on P0. On the compact outer strip S=Q intersect
\{v\textgreater=-1/(2R)\}, h\textgreater=M+1/2. For each net the compact
closed sets \{w\_(i,alpha)\textgreater=M+1/4\} in S decrease to the
empty set. Compactness and directedness give an eventual tail empty on
the WHOLE strip. Choose such tails for the finite tuple independently.
This is a finite-intersection argument for arbitrary nets, not uniform
convergence to a discontinuous endpoint.

On these tails define functions on the ENTIRE Y\^{}an:

\begin{verbatim}
Psi_(i,alpha)=max(w_(i,alpha)-h,0) on Q, and0 outside Q,
Phi_(i,alpha)=p+Psi_(i,alpha)/R.
\end{verbatim}

The strict margin on S makes Psi identically0 on one whole fixed outer
collar. Hence these are globally continuous actual H metrics: their raw
weights in the inner chart are

\begin{verbatim}
(1/R) max(w_(i,alpha)-M-1, R v),
\end{verbatim}

and on the collar and outside they are the raw weights of p.~The
original H transition functions are unchanged. The finite AFS weights
have their unit-log branch supports; an inactive zero branch is deleted
on a sufficiently small neighborhood. The max with v preserves those
same-constant supports. These verify exactly the continuous contact
bridge of Section 10.1. Thus Phi\_(i,alpha) is genuinely H-psh in one
FIXED actual ample bundle, including all field cardinalities.

They decrease pointwise, since only the w terms vary. Their explicit
global limit is

\begin{verbatim}
Phi_i=p+(1/R) max(w_i-h,0) on Q, and p outside Q.
\end{verbatim}

It is bounded and hence genuine finite energy by fixed-ample decreasing
closure. At every point in interior(P0), the floor is inactive even at
each finite stage:
w\_(i,alpha)\textgreater=w\_i\textgreater=-M\textgreater h. Consequently

\begin{verbatim}
w_(i,alpha)=R Phi_(i,alpha)+R h0+M+1,
w_i=R Phi_i+R h0+M+1.
\end{verbatim}

The exact curvature identity is

\begin{verbatim}
dd^c w_i = R(c1(H)+dd^c Phi_i) in this local construction.
\end{verbatim}

The canonical reference term R h0 cannot be dropped. Psi/R by itself is
not the actual H-potential. The local mixed measure is R\^{}n times the
restriction of the genuine RAW H mixed measure of the Phi\_i.

\subsection{10.4. Independence and weighted bounded local
convergence}\label{independence-and-weighted-bounded-local-convergence}

At finite continuous stages the displayed raw weights are identical to
the given original local weights on interior(P0). The BE/CLD continuous
product identification therefore makes their genuine H products, with
factors R, the actual local products of those weights. Compact
continuous tests extend by zero to Y. G0 bounded decreasing-net
continuity gives the local limits.

If two independently supplied local nets have endpoint w\_i on Q, the
construction uses the SAME h,p,H and R. Their global limits Phi\_i are
equal at EVERY point, by the explicit formula, including outside Q. G0
therefore gives exactly the same local limiting measure. This proves
independence of the locally supplied nets in this generated category.
Different initial collars are compared on a third shared smaller collar:
every continuous finite-stage product agrees on its inner open, and
compact tests pass those identities to each genuine bounded limit. The
positive local Radon measures consequently glue and their restrictions
agree.

For a bounded local test w0, the same collar yields w0=R Phi0+R h0+M+1
near the tested support. The Phi0 term has genuine weighted continuity
by G0, with the test-containing pairing MINUS the zero-test pairing. The
fixed local reference term is continuous and passes by ordinary weak
convergence against compact tests. Differences of tests are split into
their terms. Thus the bounded local weighted products and their limits
are approximation-independent as well.

Rational constants are sufficient. For an existing finite real-constant
AFS net, round its finitely many constants upward with uniform error
delta and add delta. The pair index (alpha,delta) is ordered by
increasing alpha and pointwise decreasing resulting weights. Choosing
alpha above two old indices and 2delta below both old deltas gives
directedness by the non-strict inequalities already proved in114. The
resulting net decreases to the SAME endpoint on the same ordinary
neighborhood. Its finite-stage uniform error tends to0, and the actual
continuous product/finite reference construction is compatible with that
limit. More explicitly, for each real-constant stage its rational
roundings, with error sufficiently small to preserve the fixed outer
strip margin, give rational-contact GLOBAL metrics in the same H. The
max operation is1-Lipschitz, so they approach that stage's explicit Phi
uniformly on all Y, with error at most2delta/R. Fixed-ample uniform
closure and continuous product compatibility make this real-constant Phi
genuine as well. The ORIGINAL real-constant Phi net therefore decreases
to the same explicit global endpoint, and G0 applies directly to it. No
interchange of unrelated local and global approximating limits is
needed. No claim about arbitrary extended-valued varying families is
made.

\subsection{10.5. Local maxima and reference-floor nonpluripolar
compatibility}\label{local-maxima-and-reference-floor-nonpluripolar-compatibility}

One collar for a finite tuple commutes with taking its maxima: the
explicit global potential for max(w,z) is max(Phi\_w,Phi\_z), since p,h
and R are identical. G0 genuine strict-winning locality and its mixed
version therefore give local strict-winning locality for bounded
generated endpoints, on the actual Borel winning sets. No continuity of
the endpoints is needed.

For the original E1 endpoint w, the retained reference b is the local
raw weight of tau. It is finite continuous on compact frame opens inside
V. The reference floors max(w,b-j) are bounded there and belong to the
generated category: apply max to actual decreasing original section
approximants and the finite original reference weight. Their intrinsic
local products have just been constructed independently of local
approximants. Appendix B, Sections B.1--B.4 identify these products with
the restrictions of the actual ample quotient floor products.

For completeness, reference independence in this LOCAL generated
construction can be checked without inferring equality merely from equal
endpoint values on a Borel set. Let b' be another finite continuous
generated local reference and bound \textbar b-b'\textbar{} by C on Q.
For fixed j,k, choose ell larger than both by more than C and use the
bounded generated lower floor

\begin{verbatim}
z_ell=max(w,b-ell,b'-ell).
\end{verbatim}

On a smaller neighborhood of the tested compact set,

\begin{verbatim}
max(w,b-j)=max(z_ell,b-j),
max(w,b'-k)=max(z_ell,b'-k).
\end{verbatim}

The common-collar bounded strict-winning theorem compares both products
with the product of z\_ell on their respective strict winning sets. They
agree on the intersection of those sets. As j and k increase, these sets
cover exactly \{w\textgreater-infinity\}. On each fixed winning set all
later restricted products agree; the restricted measures are therefore
increasing and have the same local limit. The complement, where w is
minus infinity, is assigned zero mass in this nonpluripolar
construction. For mixed products use a common ell and the joint winning
set for the finite tuple. This argument concerns generated continuous
references; it does not assert a density theorem for an arbitrary CLD
psh reference.

Because the retained quotient is actual E1\_A, G0 integrability and the
separately reconstructed E1 pole on D identify this nonpluripolar local
limit with the endpoint ample measure restricted to
\(Y^{\rm an}\setminus D^{\rm an}\). Its original zero extension has raw
mass A\^{}n and gives P zero mass. Thus the stronger
bounded-local-approximation comparison needed for the intrinsic operator
is supplied in the precise locally generated section/contact category
associated with this datum, before any target measure.

\subsection{10.6. Every requested unbounded mixed and integrable
weighted
limit}\label{every-requested-unbounded-mixed-and-integrable-weighted-limit}

The global original associated cone is the ENTIRE proposed dominated
cone, by the whole-cone converse proved in Section 5. Its E1 space
consists of actual original metrics with genuine ample E1 quotients.
Section 10.5 constructs their local operator and reference-floor
nonpluripolar extension. For any decreasing net of such metrics with E1
limits, the all-point quotient order gives genuine decreasing quotient
nets. G0 unbounded mixed continuity tests them by every original
continuous function after pullback. This supplies the original unbounded
mixed weak limit, not only the bounded local limit of Section 10.4.

An integrable actual original quotient test equals its unique quotient
on V. Every output gives P zero mass. G0 weighted continuity therefore
gives the original weighted measures with the correct pairing sign. On a
compact local frame open, raw test weights differ from that quotient by
a fixed finite continuous reference term. That term passes by weak
convergence, so the same local raw weighted identities hold. Absolute
global integrability of sigma is not required or asserted. Arbitrary
local unbounded functions outside this associated positivity category
and unrelated classical current multiplication are not inputs to the
proposed equation theory.

The construction proves independence of admissible local nets, common
collars, compatible frames and generated continuous references. Its
bounded products and relative nonpluripolar extension equal the
original-section limits of Section 8. Thus (8.1) is the genuine local
operator, before the choice of target, and Sections 5--9 establish the
full retained-datum theory on the original carrier.

\section{11. A genuinely nonprojective test over every retained
field}\label{a-genuinely-nonprojective-test-over-every-retained-field}

The geometric domain is not empty. The following is the classical
Miyake--Oda example, printed in Fujino--Sato Example 4.4, not a new
example and not a restriction of the domain of Sections 3--9.

In \(N=\mathbb Z^3\) use \[
v_1=e_1,\ v_2=e_2,\ v_3=e_3,\ v_4=(-1,-1,-1),\quad
v_5=(0,-1,-1),\ v_6=(-1,0,-1),\ v_7=(-1,-1,0).
\] The maximal cones are
\[123,\ 125,\ 256,\ 456,\ 236,\ 367,\ 467,\ 137,\ 157,\ 457.\] Start
with the four maximal cones of the fan of \(\mathbb P^3\). For an
ordered old basis \((a,b,c)\), the two additional rays are \(a+c,b+c\).
The three cones with bases \((a,b,a+c)\), \((b,a+c,b+c)\),
\((c,a+c,b+c)\) cover the nonnegative old octant in the regions
\(x\geq z\), \(x\leq z\leq x+y\), and \(z\geq x+y\), respectively. Their
shared faces match. Apply this to the ordered bases \((v_1,v_2,v_4)\),
\((v_2,v_3,v_4)\), and \((v_3,v_1,v_4)\). This gives the ten listed
cones and a complete fan. The new determinants are \(1,-1,1\) relative
to the old unimodular basis. Every maximal chart is therefore
\(\mathbb A^3_k\), over any field.

For completeness, the primary properness theorem we use is Rohrer,
dissertation IV 1.3.11--1.3.12, whose printed proof extends a DVR-valued
character into a cone of the appropriate complete star fan. Star
completeness itself has a direct proof that avoids a misleading
origin-based argument: choose \(z\) in the relative interior of the
fixed cone \(\sigma\); along a sequence \(z+\epsilon_j x\) with
\(\epsilon_j\downarrow0\), completeness and finiteness give one cone
\(\tau\) containing infinitely many terms and \(z\). Its intersection
with \(\sigma\) contains an interior point, so \(\sigma\) is a face of
\(\tau\). Projecting gives \(x\) in \(\tau/\sigma\). Thus the star fan
is complete. The DVR extension proof then establishes properness of the
toric carrier. Separatedness follows from the fan gluing theorem in the
same primary source. The charts and their shared dense torus prove
smoothness and integrality.

Every line bundle can be represented by an invariant boundary Cartier
divisor \(\sum a_iD_i\): on the dense torus the Laurent polynomial ring
is a UFD, so the line bundle is trivial there; take its rational
trivializing section and its boundary divisor. The three invariant wall
curves for \(25\), \(36\), and \(17\) are \(\mathbb P^1\), since their
quotient fans have the two primitive opposite rays and the corresponding
charts glue by inversion. Their degrees are \[
d_{25}=a_1+a_6-a_2-a_5,\qquad
d_{36}=a_2+a_7-a_3-a_6,\qquad
d_{17}=a_3+a_5-a_1-a_7.                                \tag{11.1}
\] For example, the wall relation is \(v_1+v_6=v_2+v_5\); the local
Cartier frames have characters pairing to \(-a_i\) with the respective
rays, and their difference has transition exponent \(a_1+a_6-a_2-a_5\)
on the quotient coordinate. The rational first frame section has that
order at the second chart's origin, giving the stated sign. The other
two relations are cyclic. The three numbers in (11.1) sum to zero. An
ample line bundle would have positive degree on all three proper curves,
which is impossible. Proper plus ample implies projective (Stacks 0B5N),
so this smooth proper carrier is genuinely nonprojective over every
field in our domain.

Section 3 applies to this example, as to every smooth integral proper
\(X\) and every chosen big actual rational \(L\). No toric invariance or
support restriction is used in the metric or pairing construction.

\section{Appendix A. Compatible valuative centers and local
factoriality}\label{appendix-a.-compatible-valuative-centers-and-local-factoriality}

\subsection{A.1. The precise valuative
premise}\label{a.1.-the-precise-valuative-premise}

Use Stacks tag0BX5, current Lemma29.43.1: a finite-type quasi-separated
scheme morphism is proper exactly when every valuation-ring diagram has
a unique extension. The printed proof is a formal consequence of the
scheme universal-closedness and separatedness criteria. It is not a
separately printed long proof. The scheme version suffices here; no
algebraic-space extension of the fraction field is involved.

The actual primary bodies additionally opened and read were:

\begin{itemize}
\tightlist
\item
  Section26.20, tag01KA, including Lemma26.20.5 and Proposition26.20.6.
  The former's printed proof constructs the map from the closed-point
  local ring to the given valuation ring by maximality under domination.
\item
  Lemma26.22.1, tag01KZ, with its entire printed proof: the equalizer of
  two extensions is closed and contains the generic point of the
  valuation-ring spectrum, so its defining ideal is zero in a domain.
\end{itemize}

These are the pertinent existence and uniqueness mechanisms. Earlier
topological and algebraic cross-references were not all read
recursively. The current page numbering is recorded as retrieved; old
book PDFs may have different lemma numbers. The permanent tags fix the
statements.

\subsection{A.2. Center of every analytic point on an arbitrary proper
original
X}\label{a.2.-center-of-every-analytic-point-on-an-arbitrary-proper-original-x}

Let X be a proper finite-type k-scheme, with k trivially valued, and x
in X\^{}an. Its support and completed residue field define an actual
scheme morphism Spec H(x) -\textgreater{} X over k. Put
R\_x=H(x)\^{}circ, the subring of elements of absolute value at most1.
For any nonzero element a, either a or a\^{}(-1) lies in R\_x; hence
R\_x is a valuation ring with fraction field H(x). The trivial norm on k
makes k a subring of R\_x.

Properness supplies a unique map Spec R\_x -\textgreater{} X extending
the given generic map. Its closed-point image is the center c\_X(x).
This applies to every analytic point, including points with nongeneric
support, and to every cardinality of k. No projectivity or rank-one
discreteness condition enters the valuation-ring criterion.

Let f:Y-\textgreater X be a k-morphism between proper finite-type
schemes and y in Y\^{}an, with x=f\^{}an(y). The analytic functor gives
an isometric embedding H(x)-\textgreater H(y). Thus R\_x maps locally to
R\_y: the intersection with R\_x of the maximal ideal
\{a:\textbar a\textbar\textless1\} of R\_y is the maximal ideal of R\_x.
Compose Spec R\_y -\textgreater{} Y with f.~It extends the same generic
map as Spec R\_y -\textgreater{} Spec R\_x -\textgreater{} X, so
uniqueness on X gives equality. In particular

\begin{verbatim}
f(c_Y(y))=c_X(f^an(y)).
\end{verbatim}

The identical proof works for any isometric valued-field extension and
an analytic point above x. It uses the actual residue-field map, not a
choice of a lift. It proves the compatible-center assertion needed in
original canonical metric pullback and section descent.

\subsection{A.3. Literal canonical metric
functoriality}\label{a.3.-literal-canonical-metric-functoriality}

For an actual line bundle L on proper X, choose a local algebraic frame
e at c\_X(x). The canonical trivial norm assigns norm1 to that frame at
x; a vector a e has norm \textbar a\textbar. This does not depend on the
frame: their ratio is a unit of O\_(X,c\_X(x)), whose image and inverse
both lie in R\_x, so its absolute value is1. A local regular coefficient
has norm at most1 for the same local-ring-map reason.

The compatible-center equality makes f\^{}\emph{e a frame at c\_Y(y). It
follows directly that the canonical metric on f\^{}}L is the pullback of
this canonical metric on L. Tensor products and divisible rational
powers are compatible by multiplying or taking the appropriate power of
the coefficient norms. For an actual section s of a cleared power, its
logarithmic norm therefore satisfies

\begin{verbatim}
log||f^*s||(y)=log||s||(f^an(y)),
\end{verbatim}

including the value minus infinity at a zero. This establishes the
actual metric/section pullback used in109/113 on nonprojective X. It
does not assert that the canonical metric of an arbitrary line bundle is
globally continuous or semipositive. Continuity of the specified
retained reference follows instead from its actual finite FS
representation and the compact quotient argument of110.

For a globally generating family of a cleared actual line bundle, all
its normalized section norms are at most1 at each center. Some member is
a unit coefficient in the local frame: otherwise the family would not
generate its stalk. Its norm is1. Hence the maximum of logarithmic norms
is identically0, also on a proper nonprojective carrier. This is the
generator multiplication identity used in114.

\subsection{A.4. Local factoriality source
audit}\label{a.4.-local-factoriality-source-audit}

The actual full printed statement and proof of Stacks tag0AG0, current
Lemma15.123.2, were opened and read. A regular local ring is a UFD. The
proof inducts on dimension, chooses x outside m\^{}2, makes x prime
using the regular quotient, and principalizes each height-one prime
after localization using Pic(R\_x)=0. The final irreducible factor
remains prime before localization because x is prime.

The immediately preceding tag0AFZ, Lemma15.123.1, and its entire printed
proof were also opened and read. It proves Pic(R\_f)=0 using a finite
free resolution and its determinant. Its earlier finite-free resolution,
determinant and elementary factorization cross-references are published
algebraic background; their complete bodies were not all read in this
continuation. This source audit therefore supplies the pertinent proof
units without claiming a recursive new proof of all commutative algebra.

On the smooth varieties in the retained domain, local rings are regular.
The perfect-field regular/smooth equivalence was separately bound to
Stacks0B8X in112. Thus local factoriality is applied precisely where
needed for Cartier boundary/prime-divisor statements. Normal section
division is a separate assertion and must still use its own literal
normality/division mechanism; the UFD citation does not replace that
proof or the genuine fixed-divisor cone construction.

\section{Appendix B. Original-section local products and reference
floors}\label{appendix-b.-original-section-local-products-and-reference-floors}

\subsection{B.1. The local finite stage precedes every
target}\label{b.1.-the-local-finite-stage-precedes-every-target}

Write sigma=log\textbar\textbar s\_D\textbar\textbar{} and
F(psi)=sigma+psi after pullback. Let psi\_i be finite real-valued actual
A-FS potentials. Multiplying their cleared sections by s\_(mD) and
descending them via pi\_*O\_Y=O\_X gives actual original L sections.
Their FS metrics u\_i=F(psi\_i) have the prescribed common reference
poles, and have finite continuous local weights on
\(V=X^{\rm an}\setminus P^{\rm an}\). These are actual original metrics
before any measure target.

On an algebraic frame neighborhood within V, write s\_D=f\_D e\_D and
choose an A frame e\_A whose tensor with e\_D is the pulled original L
frame. In the sign convention of105/109 the raw original local weight is

\begin{verbatim}
w_i=psi_i+h_A+log|f_D|.
\end{verbatim}

Here f\_D has no zero on the neighborhood, and both it and its inverse
are analytic. Its logarithm has zero Hessian. BE7.14 and8.3-8.5,
including the trivial-valued case, and CLD18.1.2/8.8.5 establish the
finite continuous local mixed product and the identity

\begin{verbatim}
product dd^c(w_i) = product c1(A,psi_i) on V.
\end{verbatim}

Use rational clearing and multilinearity for rational bundles. This is a
local frame/product equality for actual original section metrics, not a
target-selected or inverse-equation definition. BE7.8, whose scope is
nontrivial valuation, supplies no premise here.

\subsection{B.2. Arbitrary global type-preserving section
approximants}\label{b.2.-arbitrary-global-type-preserving-section-approximants}

There is a useful larger class than one selected upstairs net. Suppose
an actual original fixed-L FS metric a satisfies a\textless=tau+C, where
tau is the retained original reference and C is finite. At a divisorial
valuation v\_E, scale by r\textgreater0. In a cleared representation of
a, each section's logarithmic slope must be at most sigma(v\_E):
otherwise its branch beats r sigma(v\_E)+C for all sufficiently large r.
Thus each section vanishes to at least its cleared degree times the
coefficient of E in D. Normal section division gives an actual fixed-A
FS quotient after pullback and clearing rational powers.

This argument permits arbitrary real coefficients and includes every
component of D. It does not subtract two infinite endpoint values; it
divides actual sections first. The quotient metric psi\_a is genuine,
and the whole-cone theorem identifies it uniquely at all points.

Consequently, any decreasing fixed-type original FS approximating net
a\_alpha with a\_alpha\textless=tau+C and a\_alpha down to u=F(psi) has
genuine fixed-A quotients psi\_alpha. Original order off P gives
quotient order on divisorial points, then genuine ample all-point order.
Hence these quotients decrease. Their limit has the same divisorial
values as psi and is generically finite, so genuine ample decreasing
closure and divisorial determination give psi\_alpha down to psi
everywhere.

If psi is globally bounded, every psi\_alpha is bounded below by psi and
bounded above by its finite generic maximum. Every displayed fixed-A FS
approximant is therefore real-valued everywhere: a common zero of its
sections would contradict that lower bound. Thus the local finite
products of Appendix B.1 apply to every such net, not only the special
product net selected in113. For a general E1 endpoint, use bounded
reference floors first as in Appendix B.4.

\subsection{B.3. Genuine bounded decreasing limits on original frame
opens}\label{b.3.-genuine-bounded-decreasing-limits-on-original-frame-opens}

Take bounded genuine quotients psi\_1,\ldots,psi\_n.~Independently
choose decreasing finite fixed-A FS approximants to them. Appendix B.1
gives actual original local mixed products at each finite stage. Let K
be a compact subset of an original frame neighborhood contained in V,
and let g be any continuous test with support in K. Pull it through the
analytic isomorphism and extend it by zero to Y\^{}an. This is a
continuous global test: its support is a compact subset of the open
complement of D, inside compact Hausdorff Y\^{}an.

The genuine ample weak decreasing-net theorem applies to this test. It
proves convergence of the actual local products, and identifies their
limit on the frame neighborhood with the restriction of the genuine
ample mixed measure. Positivity makes these compatible Radon local
limits. Appendices B.2 and B.3 show their independence of ALL admissible
global original-section approximating nets in the retained cone.

Changing the original algebraic frame subtracts the logarithm of its
analytic unit, as in Appendix C. The finite-stage identity is unchanged,
so the weak local limits are unchanged. Restrictions to smaller frame
neighborhoods are compatible for the same test-function reason. Thus
this is an actual original-section local construction for bounded
endpoints that may be discontinuous. Uniform convergence to a
discontinuous endpoint has not been assumed.

This assertion proves more than continuous-stage compatibility. Section
10 separately constructs and compares independently supplied local nets
in the associated generated section/contact category. No assertion about
arbitrary unrelated classical currents is inferred.

\subsection{B.4. Unbounded original reference floors and nonpluripolar
extension}\label{b.4.-unbounded-original-reference-floors-and-nonpluripolar-extension}

For u=F(psi), psi in genuine E1\_A, use original metrics

\begin{verbatim}
u_j=max(u,tau-j)=F(max(psi,-j)).
\end{verbatim}

Each quotient floor is bounded genuine, so Appendix B.3 constructs its
actual local product. On V strict winning u\textgreater tau-j is exactly
psi\textgreater-j. G0's genuine ample strict-winning locality identifies
the restricted measure with the same restriction of the endpoint ample
measure. Endpoint integrability of psi makes its minus-infinity set have
zero mass. The separately proved actual ample E1 pole on D makes D have
zero mass for every finite-energy mixed output. Therefore the
strict-winning floor measures, restricted off D and read on V, increase
to the endpoint local measure and have exact eventual mass limit
A\^{}n.~Their zero extension across P is nonpluripolar in this specified
reference-floor construction and has that exact raw mass.

For mixed outputs, apply the same restriction with the common strict
winning set for the finitely many floors. Its complement shrinks to the
union of the endpoint pole sets and D, all of zero mixed mass by G0
mixed integrability and the polar test. Equivalently one can polarize
the genuine ample products; coefficients of a polarization identity are
signed and must not all be called positive.

For decreasing original finite-energy endpoints, their quotients
decrease by the genuine all-point cone/order theorem. G0 unbounded mixed
weak continuity applies to every original continuous test after
pullback. An integrable genuine original quotient test equals its unique
quotient on V and all involved measures give P zero mass. G0's correctly
signed weighted pairing identity therefore gives the corresponding
original weighted weak limits. A difference of two decreasing tests is
handled term by term; it need not decrease. No integrability of the
absolute reference sigma is asserted.

These conclusions construct an actual original operator from local
products of original sections and their reference floors before a target
is chosen. Its identification with the retained ample measure is a
conclusion of the construction, not its sole definition.

\section{Appendix C. Continuous original references and signed frame
compatibility}\label{appendix-c.-continuous-original-references-and-signed-frame-compatibility}

Retain any algebraically closed trivially valued characteristic-zero
field k, a smooth integral proper positive-dimensional original X, and
the retained projective modification pi:Y-\textgreater X, where Y is
smooth integral projective. The exact actual-ample G0 applies only on
actual ample bundles on Y. Let B be ANY actual rational line bundle on
X; the continuity argument does not require B to be big. Its canonical
metric is defined by proper valuation-ring centers as in115. Tensor,
rational-power and pullback compatibility are proved there from unique
valuative extensions. Appendix A proves these new assertions; no
nonample theorem is granted by that construction.

\subsection{C.1. Make continuity upstairs explicit by two ample
bundles}\label{c.1.-make-continuity-upstairs-explicit-by-two-ample-bundles}

On projective Y choose an actual rational ample H sufficiently large
that pi\^{}\emph{B+H is actual ample. Indeed the numerical ample cone is
open and c1(H0) is an interior point for any chosen actual ample H0. For
a sufficiently large positive integer m,
c1(H0)+m\textsuperscript{(-1)c1(pi}}B) is ample. Therefore the actual
rational bundle H=mH0 and pi\^{}*B+H are ample by numerical ampleness.
This uses only the same projective numerical ampleness premise audited
in112, including numerically trivial twists. It does not assert an ample
bundle on nonprojective X.

For either actual ample rational bundle C on Y, clear powers and choose
a finite globally generating family of sections of qC. At each analytic
point y, all their canonical norms are at most1 and at least one is
exactly1: coefficients are regular at the proper valuation-ring center,
and a generating coefficient is a local unit. Thus the section norm
maximum is identically1. In a local algebraic frame e\_C, if the
generating sections are s\_j=f\_j e\_C\^{}q, then

\begin{verbatim}
-log||e_C||_can = (1/q) max_j log|f_j|.
\end{verbatim}

The right side is finite continuous on the analytification of the frame
open. The f\_j have no common zero there, so their finite maximum is
positive before its logarithm. This proves continuity of the actual
canonical metric on C, including points with nongeneric support. Only
actual ample global generation is used.

Apply this separately to H and pi\^{}*B+H. Tensor compatibility of the
canonical metrics gives, in compatible local frames,

\begin{verbatim}
h_(pi^*B,can)=h_(pi^*B+H,can)-h_(H,can).
\end{verbatim}

The two terms are finite continuous on their common frame opens.
Consequently the canonical metric on the actual nonample pi\^{}*B is
finite continuous in every local frame. Rational powers are handled
after clearing and dividing the resulting logarithms. This proves the
particular arbitrary-bundle continuity claim needed for descent instead
of assuming it from the definition of centers.

\subsection{C.2. Closed analytic quotient and all-point
descent}\label{c.2.-closed-analytic-quotient-and-all-point-descent}

The analytic map pi\^{}an is continuous by functoriality. Y\^{}an is
compact Hausdorff in G0's projective domain. The original X\^{}an is
Hausdorff because X is separated; the pertinent actual Berkovich1993
algebraic analytification proposition is2.6.7, not a properness claim
imported from a projective-only theorem. Surjectivity of the proper
birational scheme map gives analytic surjectivity, as in the pertinent
complete printed proof of2.6.8 preserved in110.

One can see the pointwise surjectivity directly: for x in X\^{}an the
fiber of Y over Spec H(x) is nonempty and of finite type. A closed point
of that fiber has residue field finite over the complete valued field
H(x). Extend its absolute value to that finite field; it defines an
analytic point upstairs over x. This is an existence proof at each
point, with no section of pi and no choice of a measure lift.

A continuous map from compact Y\^{}an onto Hausdorff X\^{}an is closed
and quotient. In particular X\^{}an is compact. Its restriction over ANY
ordinary open U of X\^{}an is still closed and quotient: a relatively
closed F in pi\^{}(-1)(U) is F0 intersect pi\^{}(-1)(U) for closed F0 in
Y\^{}an, and its image is pi(F0) intersect U, relatively closed. This
verifies the restriction step before applying it to a local original
frame.

For an actual B-frame e on an algebraic original open U, proper
valuation-center compatibility115 gives at EVERY y over x in U\^{}an

\begin{verbatim}
||pi^*e(y)||_can = ||e(x)||_can.
\end{verbatim}

The upstairs norm is the pullback canonical metric proved continuous in
Appendix C.1. It is a finite positive continuous scalar on
pi\textsuperscript{(-1)(U}an). It is fiber-constant by the displayed
identity, not merely on the isomorphism locus or divisorial points. The
restricted quotient map therefore descends it to a finite positive
continuous norm of e on U\^{}an. Hence the canonical reference metric on
the original B is continuous, with all original transition functions and
rational powers. This completes the continuity step left deliberately
unasserted in115.

For every finite original section family, its maximum norm is now
continuous and nonnegative. Taking its logarithm gives an actual
extended-continuous FS weight, with value minus infinity precisely at
the common section zero locus. Multiplying genuine ample section
approximants by s\_D and descending sections therefore gives actual
original metrics with continuous local weights off the retained pole
set. No value on a collapsed fiber is selected after a target.

\subsection{C.3. Signed local frame transformation and arbitrary local
nets}\label{c.3.-signed-local-frame-transformation-and-arbitrary-local-nets}

Use raw metric weights w\_e=-log\textbar\textbar e\textbar\textbar. If a
new actual algebraic local frame is e'=g e for an invertible regular
function g, then

\begin{verbatim}
w_(e')=w_e-log|g|.
\end{verbatim}

The minus sign follows from the norm of the SAME metrized line bundle.
The CLD unit-log premise says dd\^{}c log\textbar g\textbar=0; both
signs are harmonic. On every compact frame overlap the term
log\textbar g\textbar{} is finite continuous.

Take ANY admissible decreasing local generated FS/contact approximation
net w\_alpha in the common-neighborhood category of117. Replace every
member by w\_alpha-log\textbar g\textbar{} on the overlap. This is still
in the same local category: actual affine-FS branches multiply by powers
of g\^{}-1 after clearing denominators, and contact supports change by
the same harmonic unit-log term and the SAME additive constant. The new
net decreases to exactly w-log\textbar g\textbar{} at every point.

At each finite continuous stage the actual local mixed products in these
two frames agree, since every curvature factor is unchanged. For a
compactly supported continuous test on the overlap,117's common
actual-ample collar supplies the weak limit in both frames. Passing the
finite-stage identity to those limits proves equality of their bounded
intrinsic local products. Independence of the generated nets in117 then
makes this true for arbitrary independently chosen nets in the second
frame, not just this matched one. Restrictions commute with compact
tests, so the local positive Radon measures glue.

The same reasoning works for a finite tuple of metrized actual bundles
with individual unit frame terms. Generated reference floors transform
as metric maxima under the SAME unit term. Bounded strict-winning
locality and reference-floor independence in117 then preserve the
unbounded nonpluripolar local products under the frame change. For a
weighted local raw test the fixed harmonic term is continuous on the
test support and passes by ordinary weak convergence; its potential term
passes by the genuine weighted theorem with the required pairing sign.
This verifies frame compatibility for the actual original operator
before any target is selected.

\section{Appendix D. Exact finite PL extensions and the genuine
contact-to-ample
bridge}\label{appendix-d.-exact-finite-pl-extensions-and-the-genuine-contact-to-ample-bridge}

Throughout this appendix, \(k\) is algebraically closed, trivially
valued and of characteristic zero. Except where smoothness is expressly
required, \(X\) is reduced irreducible projective. \(H\) and \(L\)
denote actual ample rational line bundles, not just numerical classes.
Rational PL always means rational slopes and constants. Analytic contact
weights mean the following binding hypothesis: on ordinary point
neighborhoods the continuous weight is a finite continuous
rational-affine selection in finitely many unit logarithms, with an
affine local support below it and equal to it at the tested image point.
The support inequality and equality have the \textbf{same constant}.
Finite-floor logarithmic maxima are allowed by deleting inactive
zero-coordinate branches before using unit charts. Local decreasing
approximants are supplied on one common neighborhood for each finite
tuple. These restrictions are substantive: arbitrary positive currents
and arbitrary global weak rational approximations are not asserted to
satisfy them.

\subsection{D.2. Compact algebraic extensions on the original
variety}\label{d.2.-compact-algebraic-extensions-on-the-original-variety}

\subsubsection{D.2.1. Localization and bounded
logarithms}\label{d.2.1.-localization-and-bounded-logarithms}

Fix a very ample \(H\). If \(s\) is a nonzero section, then, without
normality,
\[\Gamma(D(s),\mathcal O_X)=\varinjlim_m H^0(X,mH)/s^m. \tag{D.2.1}\] On
a finite affine trivializing cover this is module localization; the
finitely many matching equalizers commute with filtered localization.
This proves the assertion for the original structure sheaf. Thus
inverting finitely many regular functions \(\tau_i/s^{d_i}\) gives the
literal chart
\[D(a)=D(s)\cap\bigcap_iD(\tau_i),\qquad a=s\prod_i\tau_i. \tag{D.2.2}\]
The extra factor \(s\) retains the old boundary. Every regular function
on this chart is \(F/a^m\) for a global section of the appropriate power
of \(H\).

For a global section \(F\), let \(\mathfrak b_F\) be the coherent ideal
locally generated by its coefficient in a center frame. Its logarithmic
norm for the canonical trivial metric is \(\log|\mathfrak b_F|\). At
each analytic point this is independent of the center frame, since
changes of center frame are units in its valuation ring. This also
follows from BJ 1.24 and the subsequent definition of section evaluation
(S1 LF883--900). Therefore
\[A_{F,N}=\max(\log|\mathfrak b_F|,-N)\in\operatorname{PL}(X) \tag{D.2.3}\]
for positive rational \(N\), by the flag ideal \(\mathfrak b_F+(t^N)\)
after clearing denominators. If \(f=F/a^m\) is a unit on a compact set
\(W\subset D(a)^{\rm an}\), choose rational clipping levels below the
minima of both section logarithms on \(W\). Then
\[\log|f|=A_{F,N}-m A_{a,N'}\quad\hbox{on a neighborhood of }W. \tag{D.2.4}\]
The equality is exact. A uniform approximation to a logarithm is not
being substituted.

Analytic units can be treated on smaller compact rational domains as
well. In the algebraic reduction of published Ducros 3.2 (S9
LF2087--2127), a rational domain is localized at a nonvanishing
denominator \(h\), and the functions \(b/h^r\) in the original ring are
dense in its affinoid algebra. For a unit \(f\) choose an approximant
\(g\) with \[\|f-g\|_W<\min_W|f|.\] The ultrametric inequality gives
\(|g|=|f|\) at \textbf{every} point of \(W\). A finite list is treated
simultaneously. The original domain and principal localization are
retained. This uses the precise rational-domain density mechanism in
that proof, including its finite rational refinements; density of
polynomials on an arbitrary analytic open is not claimed. Norm equality
remains true after complete field extension because the same uniform
inequality and lower bound pull back.

\subsubsection{D.2.2. A finite selected weight has an actual PL
extension}\label{d.2.2.-a-finite-selected-weight-has-an-actual-pl-extension}

Suppose a finite continuous rational-affine selection is supplied near a
compact set and its used coordinates are unit logarithms there. Choose
compact rational neighborhoods inside its original charts with rational
upper/lower bounds; their interiors cover the compact set. They are
\(\Gamma\)-strict with \(\Gamma=\exp(\mathbb Q)\), even over the trivial
field. Ducros 3.2 applies to these \textbf{compact}, \(\Gamma\)-strict
spaces and their \textbf{invertible} functions. Their logarithmic images
are finite rational polytope complexes of dimension at most \(\dim X\).
No dimension equality is needed.

Refine the images by the finitely many hyperplanes where the selected
affine expressions agree. On a relative open convex cell the distinct
expressions have a fixed order. Continuity forces the selected
expression class to be locally constant, hence constant on that cell.
Continuity supplies its affine value on the closed cell. The result is a
rational PL function on a finite compact rational polyhedral complex
\(T\).

Here is an elementary extension, so no unrestricted tropical selection
theorem is hidden. Put \(T\) in the interior of a rational box. Cut that
box by all rational hyperplanes defining its cells and the expression
domains, and barycentrically subdivide the resulting finite rational
polytopes. This triangulates the box with \(T\) as a subcomplex and each
affine piece as a union of simplices. Assign the prescribed rational
value to each vertex of \(T\), assign zero to the box-boundary vertices,
and assign arbitrary rational values to the other vertices. Affine
interpolation on every simplex, followed by zero outside the box, gives
a continuous finite rational PL function \(E\) on the whole Euclidean
space, equal to the selection on \(T\).

Such \(E\) is a difference of finite convex rational PL functions. To
verify this directly, take all rational supporting hyperplanes \(h_j=0\)
of its affine cells. Across any wall, continuity makes the gradient jump
normal to the wall. There are finitely many jumps. Choose positive
rational \(C\) so adding \(C\sum_j|h_j|\) makes every oriented normal
jump nonnegative. Restriction to each generic line is then convex;
continuity handles lines through intersections. Hence \(E+C\sum_j|h_j|\)
is convex and is the maximum of its finitely many rational affine
pieces. The penalty is also a finite convex maximum. Thus \(E\) is a
finite expression using rational affine combinations, maxima and
subtraction.

Replace its unit-log coordinates by the actual global PL functions
(D.2.4). BJ's rational PL space is a vector lattice (S1 LF962--1008), so
this expression is a global algebraic PL function agreeing exactly with
the local selected weight near the chosen compact set. This construction
does not require a positive extension. It also does not identify PL on a
polytope with PL on the original bundle without constructing the
displayed algebraic functions.

\subsubsection{D.2.3. Support-controlled gluing without multiplying PL
functions}\label{d.2.3.-support-controlled-gluing-without-multiplying-pl-functions}

For compact \(K\subset U\subset X^{\rm an}\), choose a continuous bump
\(b\), equal to one on a neighborhood of \(K\), with compact support in
\(U\). Approximate it by a global rational algebraic PL \(p\) with error
less than \(1/8\), using BJ 2.2 and its lattice proof. Put
\[\chi=\min(1,\max(0,2p-1/2)). \tag{D.2.5}\] Then \(0\le\chi\le1\), it
is one near \(K\), and
\[\operatorname{Supp}\chi\subset\{p\ge1/4\}\Subset U.\] Outside \(U\),
\(p<1/8\), and the compact threshold set misses the complement of \(U\).
This proves controlled support, rather than simply extending a bump by
zero.

If finitely many local extensions \(E_i\) agree with one continuous
function \(f\) on opens containing \(\operatorname{Supp}\chi_i\), and
some \(\chi_i=1\) near every point of \(K\), choose rational \(C\)
larger than all necessary compact bounds. Then
\[E=\max_i\bigl(E_i-C(1-\chi_i)\bigr) \tag{D.2.6}\] equals \(f\) near
\(K\). Indeed a positive \(\chi_i\) puts the point in its agreement
chart, where its expression is at most \(f\); for \(\chi_i=0\) the
choice of \(C\) suppresses the bounded global \(E_i\); a winning
\(\chi_i=1\) gives equality. Only the tested compact neighborhood is
used when \(f\) is local. For a global \(f\) the same argument applies
on all of compact \(X^{\rm an}\).

Choose another cutoff \(\eta\) equal to one near \(K\) with support
compact in \(U\), and rational \(B>\sup_{\operatorname{Supp}\eta}|E|\).
The compact signed extension
\[\widetilde f=\max(-B\eta,\min(B\eta,E)) \tag{D.2.7}\] equals \(f\)
near \(K\), vanishes on a neighborhood of \(X^{\rm an}\setminus U\), and
satisfies \(|\widetilde f|\le B\eta\). Multiplication of two variable PL
functions has \textbf{not} been used. A product \(\eta E\) need not be
PL.

Given a finite open cover, choose compact shrinkings \(K_i\subset U_i\)
covering \(X^{\rm an}\) and cutoffs \(\chi_i\) as above. The actual
rational PL partition is
\[\rho_1=\chi_1,\qquad \rho_i=\min(\chi_i,1-\sum_{j<i}\rho_j). \tag{D.2.8}\]
Induction gives \(\sum_{j\le i}\rho_j=\min(1,\sum_{j\le i}\chi_j)\).
Thus every \(\rho_i\ge0\), its support is compact in \(U_i\), and
\(\sum_i\rho_i=1\). No variable denominator occurs.

\subsubsection{D.2.4. Exact differences in one actual ample
class}\label{d.2.4.-exact-differences-in-one-actual-ample-class}

For any finite list \(a_i\in\operatorname{PL}(X)\), BJ 2.25--2.26 gives
\(a_i=a_i^+-a_i^-\) with \(a_i^\pm=c_i^\pm q_i^\pm\),
\(c_i^\pm\in\mathbb Q_{\ge0}\) and \(q_i^\pm\in H(H)\). Choose a common
rational \(R\ge c_i^\pm\) and set
\[q_{i,R}^\pm=c_i^\pm q_i^\pm+(R-c_i^\pm)0\in H(RH). \tag{D.2.9}\] The
zero is the actual canonical FS metric on \(H\). Tensor products of the
generating sections prove the asserted FS membership; denominators are
cleared once for the finite list. Thus \(a_i=q_{i,R}^+-q_{i,R}^-\)
exactly. In original frames \(e_\alpha=g_{\alpha\beta}e_\beta\) of
\(H\), both terms have weights \(Rh_{0,\alpha}+q_{i,R}^\pm\) and
transition \(g_{\alpha\beta}^{NR}\) on the cleared tensor power. Their
difference is \(a_i\). This is the same algebraic bundle, same cocycle,
same reference and same class. Positivity of an abstract numerical class
alone is not used.

\subsection{D.3. Audit and reconstruction of the continuous contact
metric
bridge}\label{d.3.-audit-and-reconstruction-of-the-continuous-contact-metric-bridge}

Let a continuous rational contact metric be \textbf{already glued} on an
ample rational \(L\). In a standard very ample frame \(e=s_j\) of
\(qL\), its canonical reference weight is
\[h_{0,j}=q^{-1}\max_l\log|s_l/s_j|,\] with the \(l=j\) term giving a
zero floor. Its actual potential is \(\phi=b_j-h_{0,j}\), independent of
the frame. The finite-floor zero reductions and Appendix D.2 extend its
local rational selected weights exactly. Formula (D.2.6), on the finite
compact cover of \(X^{\rm an}\), proves \[\phi\in\operatorname{PL}(X)\]
on the original variety and original \(L\).

BJ 2.7/2.31 now supplies an actual algebraic test-configuration line
bundle \(\mathcal L\) with generic bundle \(L\) and model function
\(\phi\). Take relative integral closure inside that prescribed generic
fiber, not absolute normalization. Complete at the test parameter after
clearing rational tensor weights over \(K=k((t))\), \(|t|=e^{-1}\). The
model metric is exactly the pullback of the original metric: for a flag
ideal its norm is the maximum of the norms of the \textbf{separate}
ideal generators \(a_\lambda t^{-\lambda}\); differences preserve
equality. This is an equality at all \(K\)-analytic points, not only the
Gauss section. The model identification uses the original \(L\) and its
original cocycle.

This is a new resolution of the algebraization dependency in the
supplied bridge. GM 2.4 algebraizes a formal blowup ideal, while GM 2.6
cites the separate GK17 8.13 for arbitrary formal line-bundle
algebraization (S4 LF249--295). That latter full proof is not in the
authorized caches. Here no arbitrary formal line bundle is algebraized:
the algebraic test configuration was constructed \textbf{first} by BJ's
flag-ideal representation. Proper formal-model refinements in the
following primary criteria remain background theorems with their precise
applications; a title is not presented as their proof.

Over a completed algebraic closure of \(K\), pull back a contact support
through any smooth analytic curve map into an ordinary weight chart. It
becomes \[h_x=\sum_j a_j\log|f_j|+C,\] harmonic by Thuillier 2.3.21. If
\(b-h\) has a local maximum \(M\) for any harmonic \(h\), then
\(h_x-h\le b-h\le M\) with equality at \(x\). The harmonic maximum
principle forces \(h_x-h=M\) nearby and then \(b-h=M\) nearby. Thuillier
3.1.12 gives subharmonicity. This proves the curve test from supplied
local expressions; it does not pull back arbitrary positive currents.
For decreasing local approximants, restrict \textbf{each} one first and
use Thuillier 3.1.8(ii) afterwards. A bounded floor excludes the
identically \(-\infty\) alternative.

The completed model metric is already rational PL on the original
bundle. GR 5.1/5.4, whose full proofs are pinned at S8 LF803--926,
therefore give semipositivity on each ordinary point neighborhood. GR's
proof lifts the generic reduction point of an entire proper residual
curve, retaining the interior condition; its local curve-lifting theorem
is pinned at LF489--800. GM 3.11 tests every proper residual curve of
the represented global model, using its generic reduction point, a
neighborhood and a common dominating model. The closure of that point's
reduction remains the \textbf{entire} proper curve (S4 LF754--782). GM
3.5 and 3.3 then test nefness on the already supplied algebraic model
and detect it before field extension. There is no positivity assertion
from a Laurent G-cover alone. GR Remark 1.8 leaves the unrestricted
current converse unknown.

BE 7.8 now applies over the nontrivial \(K\) to the ample \(L_K\). It
interpolates the nef model with an ample model of the \textbf{same}
generic bundle. The resulting semiample models give uniform pure FS
approximation. The printed \(H/A\) switch in its proof is resolved by
using one ample model consistently.

For exact descent, expand each pure FS section on \(mL_K\) in the finite
\(k\)-space \(H^0(X,mL)\): \[s_j=\sum_{r\ge r_0}t^r f_{jr}.\] At the
independent-variable Gauss lift of \(x\), the FS potential is
\[m^{-1}\max_{j,r}\{\log\|f_{jr}(x)\|_{mh_0}-r\}. \tag{D.3.1}\] Every
\(k\)-section has norm at most one in the canonical metric, by its
valuation-ring center. The left side is finite continuous on compact
\(X^{\rm an}\), so has a finite lower bound. Every sufficiently large
\(r\) term lies uniformly below that bound. The infinite maximum is
therefore \textbf{exactly} a finite cutoff. Retained coefficients
generate, since a common algebraic zero would contradict finiteness and
the same tail bound. They are sections of the original \(mL\), with
rational weights. Uniform errors survive Gauss restriction. BE 7.6 or BJ
decreasing closure yields genuine continuous PSH on the original \(L\)
over \(k\).

For a continuous target choose uniform errors \(\epsilon_j=4^{-j}\) and
shifts \(B_j=2\sum_{i\ge j}\epsilon_i\). Then shifted FS approximants
decrease, since the consecutive lower bound is
\(2\epsilon_j-\epsilon_j-\epsilon_{j+1}>0\). This proves the continuous
contact bridge without any assertion that every continuous weak metric
has global rational weak approximants.

\subsection{D.4. Genuine collars and bounded local weighted
products}\label{d.4.-genuine-collars-and-bounded-local-weighted-products}

The supplied round75 collar is reconstructed, including its reference.
On a principal chart (D.2.2), write the finite defining constraints as
\(F_j/a^m\). Enlarge \(m\) so its tensor power is very ample and append
a complete section system; its ratios generate the \textbf{original}
affine chart ring. Rational strict margins and a negative floor give an
actual FS local weight \(v\) on a fixed ample \(H\), with
\[Q=\{v\le0\}\Subset U,\qquad P=\{v\le-\epsilon\},\qquad x\in\operatorname{int}P.\]
If the tensor denominator is \(q\), the weight is
\(q^{-1}\max(q\lambda,\log|F_j/a^m|-C_j)\); the thresholds are \(C_j\),
not \(qC_j\). Generator bounds ensure compactness. Write \(v=h_0+p\) in
that frame, with \(p\) its actual global FS potential and
\(\omega=c_1(H)\).

Let \(u_i\) be a finite tuple of bounded local decreasing contact limits
on a neighborhood of \(Q\), including a bounded test if needed. Choose
rational \(M\) bounding the limits there and rational \(A\) with
\(A\epsilon>2M+2\). Put \(h=Av+M+1\). On \(P\), \(h<-M\); on the compact
outer strip \(v\ge-1/(2A)\), \(h\ge M+1/2\). The closed sets
\(\{u_{ij}\ge M+1/4\}\) on this strip decrease to the empty set. The
finite-intersection property gives one common eventual tail strictly
below \(h\) on the \textbf{whole} strip. This is not Dini's theorem for
a discontinuous limit.

Define
\[\Psi_{ij}=\max(u_{ij}-h,0)\text{ on }Q,\qquad \Psi_{ij}=0\text{ outside},\]
and \[\Phi_{ij}=p+A^{-1}\Psi_{ij}. \tag{D.4.1}\] Their actual local
metric weights on \(H\) are \(A^{-1}\max(u_{ij}-M-1,Av)\) inside, and
\(v\) on a whole reference-winning collar and outside. The contact
bridge, applied after clearing \(A\), proves that these are genuine
metrics on the \textbf{same} \(H\). Their original cocycle is unchanged.
The bounded decreasing limits \(\Phi_i\) are genuine finite energy by
global decreasing closure. On \(\operatorname{int}P\),
\[u_i=A\Phi_i+Ah_0+M+1,\qquad
dd^c u_i=A(\omega+dd^c\Phi_i). \tag{D.4.2}\] This is the exact reference
identity. \(\Psi_i/A\) alone is not a canonical-reference potential. The
local product is the restriction of \(A^n\) times the genuine
\textbf{raw} mixed measure of the \(\Phi_i\).

For continuous inputs this equals the continuous local product: BJ 7.23
invokes actual BE 8.18, whose proof and fundamental-cycle conventions
are pinned; BJ 7.42 supplies continuous locality. CLD 8.7.2 is used only
within its continuous locally-uniform premise (S3 LF10432--10449). For
bounded decreasing inputs, use BJ 7.19/7.21 on the actual global
\(\Phi_{ij}\).

Independence is checked on a third smaller shared collar. At every
continuous stage the old and smaller collar products agree on their
common inner open, because their actual weights and scaling agree there.
Pass both genuine decreasing limits against compact continuous tests. In
the shared collar different supplied sequences have the same explicit
global limit \(\max(u_i-h,0)\), zero outside. Their genuine measures
therefore agree. This proves sequence, collar and reference
independence; the local positive Radon measures glue.

For the bounded test, (D.4.2) is a genuine decreasing potential plus a
fixed continuous local reference. BJ 7.21 supplies weighted convergence
for its genuine part; ordinary weak convergence supplies the reference
part against compact cutoffs. In that source the correct identity is the
test-containing pairing \textbf{minus} the zero-test pairing (S1
LF4108--4113); the opposite signs at LF4119--4123 are a printed error.
Translate the test below zero, apply the positive weighted-measure lemma
to its negative, and subtract in the correct order. Signed tests that
are differences of such bounded weights are treated term by term; their
difference need not decrease.

\subsection{D.5. Compact signed IBP from actual global
pairings}\label{d.5.-compact-signed-ibp-from-actual-global-pairings}

Let \(a\) be a rational algebraic PL test with support \(K\Subset U\),
and let \(u_0,u_1,\ldots,u_{n-1}\) be bounded local contact limits on a
common neighborhood of \(K\). Choose one genuine collar for the finite
tuple and an additional compact cutoff equal to one near \(K\). For
\(a=q^+-q^-\) use (D.2.9) in one actual ample class \(RH\). The positive
global mixed measures obtained by adding \(q^+\) or \(q^-\) separately
need \textbf{not} be supported in \(K\). Their \textbf{difference} is:
on the open complement of \(K\), \(a=0\), so \(q^+=q^-\) and genuine
mixed locality makes their measures equal there. The same argument works
on every open where \(a\) is constant. Cancellation is retained before
any outside contribution is discarded.

For each collar, extend the local reference term in (D.4.2) to a compact
signed algebraic PL \(G_0\) agreeing with \(Ah_0+M+1\) near \(K\), using
Appendix D.2. Then \[F_0=A\Phi_0+G_0\in\widetilde{\mathcal E}^{1}\] is
an actual global signed finite-energy function and equals \(u_0\) near
\(K\). The zero-class curvature of \(F_0\) agrees there with
\(dd^cu_0\). To verify the reference cancellation, do it first for the
continuous collar approximants, using continuous local compatibility.
The difference of the reference curvature and the curvature of the
matching \(G_0\) is zero on the same open. Pass the bounded genuine
mixed limits in each positive term. Thus that cancellation still holds
with the discontinuous bounded other factors.

The symmetric global mixed pairing now gives
\[\int_U a\,dd^cu_0\wedge\bigwedge_{i=1}^{n-1}dd^cu_i
=\int_U u_0\,dd^ca\wedge\bigwedge_{i=1}^{n-1}dd^cu_i. \tag{D.5.1}\]
Precisely, both sides are the global pairing with the two zero-class
entries \((0,a)\) and \((0,F_0)\) in opposite orders, multiplied by the
other collar scaling factors. The left test vanishes outside \(K\); the
signed measure on the right vanishes there by the cancellation just
proved. Every remaining integral is of bounded functions over a compact
set against finite signed Radon measures. This proves compact IBP
without first assuming a weak/global theorem.

One can equivalently prove (D.5.1) at continuous stages and pass each
positive global mixed term using weighted convergence from Appendix D.4.
The total variations are bounded by the sum of their fixed positive raw
intersection masses. No discontinuous local limit is called locally
uniform. No difference of decreasing sequences is called decreasing.
Extension of each \(u_i\) on all of \(U\) is neither needed nor
asserted; matching near the controlled support suffices.

\section{References and verification
scope}\label{references-and-verification-scope}

The retained actual-ample theorem is the full original P16.1 theorem
stated in Section 1; its proof and separate source-aware review
accompany the P16.1 paper. The present new nonprojective construction
was reconstructed independently using Xean's native Codex path with
gpt-6.1-sol at xhigh reasoning. Review115 accepted this complete proof
packet, with 32 exact external premises and 54 source passages. The
canonical report is preserved separately. The manuscript and all
mathematical annexes appear in that frozen packet in their entirety;
this edition changes their arrangement and resolved-status wording,
while preserving their proof content. It imports none of the later
unreviewed continuations120--125.

The review does not assert reading every primary work or all recursive
background proofs. Primary statements and pertinent proof units, rather
than titles, hashes or compilation, support the applications. In
particular, BE's continuous theorem is applied only in its uniform
continuous category, and its nontrivial-valued nef theorem only over the
auxiliary \(k((t))\). G0's decreasing-net theorem supplies the
discontinuous and finite-energy passages after genuine ample membership
is established.

The permanent Stacks tags used here are
\href{https://stacks.math.columbia.edu/tag/0200}{0200},
\href{https://stacks.math.columbia.edu/tag/07QW}{07QW},
\href{https://stacks.math.columbia.edu/tag/0B8X}{0B8X},
\href{https://stacks.math.columbia.edu/tag/0AY8}{0AY8},
\href{https://stacks.math.columbia.edu/tag/0AG0}{0AG0},
\href{https://stacks.math.columbia.edu/tag/01Q3}{01Q3},
\href{https://stacks.math.columbia.edu/tag/0B5T}{0B5T},
\href{https://stacks.math.columbia.edu/tag/0BEM}{0BEM},
\href{https://stacks.math.columbia.edu/tag/0BJ8}{0BJ8},
\href{https://stacks.math.columbia.edu/tag/0BFP}{0BFP},
\href{https://stacks.math.columbia.edu/tag/0BX5}{0BX5},
\href{https://stacks.math.columbia.edu/tag/01KA}{01KA},
\href{https://stacks.math.columbia.edu/tag/01KZ}{01KZ},
\href{https://stacks.math.columbia.edu/tag/0AFZ}{0AFZ},
\href{https://stacks.math.columbia.edu/tag/01KF}{01KF} and
\href{https://stacks.math.columbia.edu/tag/0B5N}{0B5N}.

The other primary sources are:

\begin{itemize}
\tightlist
\item
  Berkovich, \emph{Étale cohomology for non-Archimedean analytic
  spaces}, 1993,
  \href{https://www.numdam.org/item/10.1007/BF02712916.pdf}{§2.6}.
\item
  Boucksom--Jonsson, \emph{Global pluripotential theory over a trivially
  valued field},
  \href{https://arxiv.org/pdf/1801.08229v4}{1801.08229v4}.
\item
  Boucksom--Eriksson, \emph{Spaces of norms, determinant of cohomology
  and Fekete points in non-Archimedean geometry},
  \href{https://arxiv.org/pdf/1805.01016v3}{1805.01016v3}.
\item
  Chambert-Loir--Ducros, \emph{Formes différentielles réelles et
  courants sur les espaces de Berkovich},
  \href{https://arxiv.org/pdf/1204.6277v2}{1204.6277v2}.
\item
  Conrad, \emph{Ostrowski's theorem and finite extensions},
  \href{https://math.stanford.edu/~conrad/676Page/handouts/ostrowski.pdf}{§6}.
\item
  Ducros, \emph{Espaces de Berkovich, polytopes, squelettes et théorie
  des modèles},
  \href{https://www.numdam.org/item/10.1142/S1793744212500077.pdf}{published
  Theorem3.2}.
\item
  Fujino--Miyamoto, \emph{Nakai--Moishezon ampleness criterion for real
  line bundles}, \href{https://arxiv.org/pdf/2101.00806v1}{2101.00806v1,
  Theorem1.4}.
\item
  Fujino--Sato,
  \href{https://www.math.kyoto-u.ac.jp/~fujino/wSato.pdf}{Example4.4}.
\item
  Gubler--Martin, \emph{On Zhang's semipositive metrics},
  \href{https://ems.press/content/serial-article-files/26466}{2019}.
\item
  Gubler--Rabinoff,
  \href{https://arxiv.org/pdf/2503.05513v1}{2503.05513v1} and
  \href{https://arxiv.org/pdf/2506.13548v1}{2506.13548v1}.
\item
  Rohrer,
  \href{https://fredrohrer.ch/wp-content/uploads/2020/06/diss.pdf}{dissertation,
  Chapter IV}.
\item
  Temkin, \emph{Desingularization of quasi-excellent schemes in
  characteristic zero},
  \href{https://arxiv.org/pdf/0904.1592v2}{0904.1592v2}.
\end{itemize}

The auxiliary toric example tests genuine nonprojectivity and does not
restrict the whole-domain theorem or the admissible measures. All
prescribed types, operators and masses are determined before a target is
selected.

\end{document}
